\subsection{对称幂的上同调}\etiquette[5.5]{XVII.subsec5.5}
\pageoriginale

本 ${\rm n^{\circ}}$ 的主要结果是定理 \hyperref[XVII.thm5.5.21]{5.5.21}; 特别地, 对适当系数 (例如 $\ZZ/n\ZZ$), 它用代数闭域 $k$ 上概形 $X$ 的固有支集上同调, 给出其对称幂的固有支集上同调, 也就是以 $X$ 的固有支集上同调来表示. 这一定理在半单纯拓扑中的类似物是众所周知的, 至少对以交换环 $A$ 为系数的上同调而言如此; 它可在半单纯上链的层面验证. 事实上, 若 $X_{\ast}$ 为半单纯集, 则从 $X_{\ast}$ 的第 $n\sisi{}{}$ 对称幂 $(X_{\ast})^n/\Sf_n$ 到 $X_{\ast}$ 的几何实现的第 $n$ 对称幂, 有一个典范连续双射, 并且它在上同调上诱导同构.

\subsubsection{}\etiquette[5.5.1]{XVII.subsubsec5.5.1}
回顾平坦模的一些性质 (归功于 \textsc{D. Lazard}), 函子 $TS^n$ 与 $T^n$ 的一些性质 (归功于 \textsc{Roby}), 以及非加性函子的导出函子的一些性质 (归功于 \textsc{Dold} 和 \textsc{Puppe}).

\paragraph{}\etiquette[5.5.1.1]{XVII.par5.5.1.1}
\emph{环 $A$ 上的每个平坦模都是有限生成自由模的滤过归纳极限} (\textsc{Lazard} \cite[I 1.2]{sga4.exp.XVII.bib.La}). 这一结果蕴含下述陈述:

\paragraph{}\etiquette[5.5.1.2]{XVII.par5.5.1.2}
函子 $\varinjlim$ 给出一个等价, 从有限生成自由 (\resp 有限生成投射) $A$-模范畴的 \iInd-对象 (\hyperrefext[I][I.sec8.1]{8.1}) 范畴到平坦 $A$-模范畴.

\paragraph{}\etiquette[5.5.1.3]{XVII.par5.5.1.3}
若\pageoriginale $A$ 与 $B$ 为两个环, 则每个将有限生成自由模映为平坦模并与滤过归纳极限可交换的函子 $T:\Mod(A)\rightarrow\Mod(B)$, 都将平坦模映为平坦模 (\textsc{Lazard} \cite[I 1.5]{sga4.exp.XVII.bib.La}).

\paragraph{}\etiquette[5.5.1.4]{XVII.par5.5.1.4}
每个 $A$-模的正合列, 若 $L''$ 平坦,

$$
0\longrightarrow L'\longrightarrow L\longrightarrow L''\longrightarrow
0
$$
就是分裂正合列的滤过归纳极限 (\textsc{Lazard} \cite[I 2.3]{sga4.exp.XVII.bib.La}).

\subsubsection{}\etiquette[5.5.2]{XVII.subsubsec5.5.2}
设 $A$ 为交换环, $M$ 为 $A$-模.

\paragraph{}\etiquette[5.5.2.1]{XVII.par5.5.2.1}
对每个整数 $n$, 以 $TS^{n}(M)$ (或 $TS^{n}_A(M)$) 表示 $M$ 的 $n$ 次对称张量（symmetric tensor）模. 以 $\Sf_n$ 表示 $n$ 次对称群, 则按定义有
$$
TS^{n}(M)=(\bigotimes^{n}M)^{\Sf_n}{\quoi}.
$$

若 $B$ 为 $A$-代数, 则 $TS^{n}(B)$ 是 $\bigotimes\limits^{n}B$ 的子代数. 若 $B$ 还交换, 则按定义 $\Spec (TS^{n}(B))$ 为 $\Spec(B)$ 在 $\Spec(A)$ 上的第 $n$ 对称幂:
$$
\Spec(TS^{n}_{A}(B))=\Sym^{n}_{\Spec(A)}(\Spec(B)\underset{\dfn}{=}(\Spec(B)/\Spec(A))^{n}/\Sf_n{\quoi}.
$$

\paragraph{}\etiquette[5.5.2.2]{XVII.par5.5.2.2}
回顾 (\textsc{Roby} \cite[I \sisi{1 pg.226}{2 p.~219}]{sga4.exp.XVII.bib.Ro}), 从 $A$-模 $M$ 到 $A$-模 $N$ 的一个\emph{多项式律}（polynomial law）$f$, 是对每个交换 $A$-代数 $A'$ 给定映射 $f_{A'}:M\otimes_{A}A'\rightarrow N\otimes_{A}A'$, 并且这些映射满足: 对每个 $A$-代数同态 $u:A'\rightarrow A''$, 图
\[
\xymatrix{
M \otimes_A A'  \ar[r]^{f_{A'}} \ar[d]^{u}& N \otimes_A A' \ar[d]^{u}\\
M \otimes_A A'' \ar[r]^{f_{A''}}& N \otimes_A A''
}
\]
均\pageoriginale 交换.

若对每个 $a'\in A'$ 恒有
\[
f_{A'}(a'x)={a'}^{n}f_{A'}(x){\quoi}.
\]
则称多项式律 $f$ \emph{为 $n$ 次齐次的} (\textsc{Roby} \cite[I \sisi{1}{8} p\sisi{g}{}.~226]{sga4.exp.XVII.bib.Ro}). 这样的多项式律也称为\emph{$n$ 次映射} (二次、三次、\ldots); 它们所成的集合记为
\[
\Hom n^{\text{\rm ique}} (M,N){\quoi}.
\]

更一般地, 若 $(M_{i})_{i\in I}$ 为一族有限个 $A$-模, $\sheafn =(n_{i})_{i\in I}$ 为一族整数, 则从模 $M_{i}$ 到 $N$ 的多重次数为 $\sheafn $ 的多重齐次多项式律, 是从 $\sisi{\sum}{\oplus} M_{i}$ 到 $N$ 的多项式律, 使对 $a_{i}\in A'$ 恒有
\[
f_{A'}((a_{i}x_{i})_{i\in I})=\prod_{i\in I} a^{n_{i}}_i\cdot
f_{A'}((x_{i})_{i\in I}){\quoi}.
\]
以 $\Hom \sheafn ^{\text{\rm ique}}((M_{i}),N)$ 表示这些多项式律所成的集合.

\paragraph{}\etiquette[5.5.2.3]{XVII.par5.5.2.3}
对每个模 $M$, 从 $A$-模到集合的协变函子
$$
N\longmapsto \Hom n^{\text{\rm ique}} (M,N)
$$
由一个 $A$-模 $\Gamma^{n}(M)$ (或 $\Gamma^{n}_{A}(M)$) 表示, 该模配有\emph{泛\pageoriginale $n$ 次映射} $\gamma^{n}:M\rightarrow
\Gamma^{n}(M)$ (\textsc{Roby} \cite[IV 1 \sisi{pg.226}{p.~265}]{sga4.exp.XVII.bib.Ro}). 给定如 \hyperref[XVII.par5.5.2.2]{5.5.2.2} 中的 $(M_{i})_{i\in I}$ 与 $\sheafn $, 协变函子
$$
N\longmapsto \Hom \sheafn ^{\text{\rm ique}} ((M_{i}),N)
$$
由模 $\bigotimes\limits_{i\in
I}\Gamma^{n_i}(M_{i})$ 表示, 该模配有\emph{泛多重齐次多项式律 (多重次数为 $\sheafn $)} $(m_{i})\mapsto
\bigotimes\limits_{i\in I}\gamma^{n_{i}}(m_{i})$.

\paragraph{}\etiquette[5.5.2.4]{XVII.par5.5.2.4}
函子 $TS^{n}$ 与 $\Gamma^{n}$ 都与滤过归纳极限可交换 (后者见 \textsc{Roby} \cite[IV \sisi{6 pg.275}{7 p.~277}]{sga4.exp.XVII.bib.Ro}). 它们将自由模映为自由模, 因而也将平坦模映为平坦模 (\hyperref[XVII.par5.5.1.2]{5.5.1.2}).

\paragraph{}\etiquette[5.5.2.5]{XVII.par5.5.2.5}
由 $m\mapsto \bigotimes\limits^{n}m$ 给出的从 $M$ 到 $TS^{n}(M)$ 的 $n^{\text{\rm ique}}$ 次映射, 确定了一个从 $\Gamma^{n}(M)$ 到 $TS^{n}(M)$ 的映射 $\varphi$, 称为典范映射, 满足
$$
\varphi(\gamma^{n}(m))=\bigotimes^{n}m{\quoi}.
$$
当 $M$ 有限自由时, 这一映射为同构 (\textsc{Roby} \cite[IV \sisi{pg.272}{5 p.~272}]{sga4.exp.XVII.bib.Ro}); 因而当 $M$ 平坦时也为同构 (\hyperref[XVII.par5.5.2.4]{5.5.2.4} 及 \hyperref[XVII.par5.5.1.1]{5.5.1.1}).

设 $M$ 为 $A$-模, $P$ 为正合列
$$
P:P_{1}\xrightarrow{d} P_{0} \xrightarrow{\epsilon} M\longrightarrow
0{\quoi}.
$$
令 $P'_{1}=P_{1}\times P_{0}$, 并设 $j_{0}$ 与 $j_{1}$ 为从 $P'_{1}$ 到 $P_{0}$ 的态射, 坐标分别为 $(d,\Id)$ 与 $(0,\Id)$. 则序列
$$
\xymatrix@C=1.75cm{
\Gamma^{n}(P'_1) \ar@<2pt>[r]^{\Gamma^{n}(j_0)} \ar@<-2pt>[r]_{\Gamma^{n}(j_1)}& \Gamma^{n}(P_0) \ar[r]& \Gamma^{n}(M) \ar[r]& 0
}
$$
正合 (\textsc{Roby} \cite[\sisi{pg.284}{IV 10 p.~284}]{sga4.exp.XVII.bib.Ro}). \emph{当 $P$ 为 $M$ 的平坦表示时},\pageoriginale 这就用函子 $TS^{n}$ 给出了函子 $\Gamma^{n}$ 的描述:
$$\etiquette[5.5.2.5.1]{XVII.eq5.5.2.5.1}
TS^{n}(P'_1)\rightrightarrows TS^{n}(P_{0})\rightarrow
\Gamma^{n}(M)\rightarrow 0{\quoi}.\leqno(5.5.2.5.1)
$$

\paragraph{}\etiquette[5.5.2.6]{XVII.par5.5.2.6}
若 $M=M'\oplus M''$, 则 $n^{\text{\rm ique}}$ 次映射
\[
m\longmapsto (\gamma^{n'}m\otimes \gamma^{n''}m)_{n'+n''=n}
\]
定义一个称为典范同构的同构 (\textsc{Roby} \cite[III \sisi{4}{9 p.~262}]{sga4.exp.XVII.bib.Ro})
$$\etiquette[5.5.2.6.1]{XVII.eq5.5.2.6.1}
\Gamma^{n}(M)\isomto \sisi{\sum}{\bigoplus}_{n'+n''=n} \Gamma^{n'}(M')\otimes \Gamma^{n''}(M''){\quoi}.\leqno(5.5.2.6.1)
$$

若 $\Sigma$ 为正合列
\[
\Sigma:0\longrightarrow M'\longrightarrow M\longrightarrow
M''\longrightarrow 0{\quoi},
\]
则模 $\Gamma^{n}(M)$ 配有典范的 (递减) 滤过及态射
$$\etiquette[5.5.2.6.2]{XVII.eq5.5.2.6.2}
\Gamma^{i}(M'')\otimes \Gamma^{n-i}(M')\longrightarrow \Gr^{i}(\Gamma^{n}(M)){\quoi}.\leqno(5.5.2.6.2)
$$
这些数据的特征为: 它们关于 $\Sigma$ 具有函子性, 且当 $\Sigma$ 分裂时满足与 \eqhyperref[XVII.eq5.5.2.6.1]{5.5.2.6.1} 的显然相容性. 当 $\Sigma$ 分裂时, 箭头 \eqhyperref[XVII.eq5.5.2.6.2]{5.5.2.6.2} 为同构, 因而当 $M''$ 平坦时也为同构 (\hyperref[XVII.par5.5.1.4]{5.5.1.4}).

\paragraph{}\etiquette[5.5.2.7]{XVII.par5.5.2.7}
函子 $\Gamma^{n}$ 与标量扩张可交换 (\textsc{Roby} \cite[III \sisi{3 pg.262}{8 p.~262}]{sga4.exp.XVII.bib.Ro}) (对于 $TS^{n}$ 的类似性质一般不成立).

\paragraph{}\etiquette[5.5.2.8]{XVII.par5.5.2.8}
若 $u:M_{1}\times M_{2}\rightarrow N$ 为双线性映射, 则 $\gamma^{n}\circ u$ 为双次数 $(n,n)$ 的双齐次多项式律, 从而确定 (\hyperref[XVII.par5.5.2.3]{5.5.2.3}) 一个满足 $u^{n}(\gamma^{n}x\otimes
\gamma^{n}y)=\gamma^{n}u(x,y)$ 的 $u^{n}:\Gamma^{n}(M_{1})\otimes
\Gamma^{n}(M_{2})\rightarrow \Gamma^{n}(N)$. 特别地,\pageoriginale 若 $B$ 为 $A$-代数, 乘法 $B\times B\rightarrow B$ 在 $\Gamma^{n}(B)$ 上诱导一个乘积. 这一乘积使 $\Gamma^{n}(B)$ 成为 $A$-代数, 并且从 $\Gamma^{n}(B)$ 到 $TS^{n}(B)$ 的典范映射是代数同态.

\subsubsection{}\etiquette[5.5.3]{XVII.subsubsec5.5.3}
设 $\Aa$ 为阿贝尔范畴. 以 $\text{\rm ssimpl}^{+}(\Aa)$ 表示 $\Aa$ 中的余半单纯对象\nde{现代术语称为\emph{余单纯}对象（cosimplicial objects）.}所成的范畴, 以 $C^{\geq 0}(\Aa)$ 表示满足在 $i<0$ 时 $K^{i}=0$ 的微分复形 $K$ 所成的范畴, 以 $K^{\geq
  0}(\Aa)$ 表示 $C^{\geq 0}(\Aa)$ 在 $K(\Aa)$ 中的像. 若 $K\in \Ob
\text{ssimpl}^{+}(\Aa)$, 则伴随的正规化复形（normalized complex）$N(K)\in \Ob C^{\geq 0}(\Aa)$ 定义为
$$\etiquette[5.5.3.1]{XVII.eq5.5.3.1}
\begin{cases}
N(K)^{n} =\text{退化算子的核的交}\\
\qquad\qquad\text{(从 $K^{n}$ 到 $K^{n-1}$)}\\
d=\sum(-1)^{i}\partial_{i}{\quoi}.
\end{cases}\leqno(5.5.3.1)
$$

若 $K$ 也表示各分量为 $K^{n}$, 微分为 $d=\sum(-1)^{i}\partial_{i}$ 的复形, 则 $N(K)$ 为 $K$ 的直和因子子复形, 且 $K/N(K)$ 零伦.

由 \textsc{Dold-Puppe} \cite{sga4.exp.XVII.bib.DP}, 函子 $N$ 为范畴等价
$$
N:\;\text{\rm ssimpl}^{+}(\Aa)\longrightarrow C^{\geq
0}(\Aa){\quoi}.
$$
这一等价及其逆都将同伦态射映为同伦态射. 此外, 对每个 $K\in \Ob\text{\rm
ssimpl}^{+}(\Aa)$, $K^{n}$ 为若干个 $N(K)^{i}$ 的直和 $(i\leq n)$; 若 $\PP$ 为 $\Aa$ 中对象的一个性质, 在有限直和与直和因子下稳定, 则 $K^{i}$ 在 $i\leq n$ 时具有性质 $\PP$, 当且仅当 $N(K)^{i}$ 在 $i\leq n$ 时具有这一性质.

\textsc{Dold} 和 \textsc{Puppe} 利用等价\pageoriginale $N$, 将每个从 $\Aa$ 到任意范畴 $\Bb$ 的函子 $T:\Aa\rightarrow \Bb$ (未必加性), 延拓为从 $C^{\geq 0}(\Aa)$ 到 $C^{\geq
0}(\Bb)$ 的函子 $C^{\geq 0}
(T)$ (或简记 $T$):
$$\etiquette[5.5.3.2]{XVII.eq5.5.3.2}
C^{\geq 0}(T)=NTN^{-1}{\quoi}.\leqno(5.5.3.2)
$$

函子 $C^{\geq 0}(T)$ 下降到商范畴后定义
$$\etiquette[5.5.3.3]{XVII.eq5.5.3.3}
T\quad\mbox{\rm 或}\quad K^{\geq 0}(T):K^{\geq
0}(\Aa)\longrightarrow K^{\geq 0}(\Bb){\quoi}.\leqno(5.5.3.3)
$$

\setcounter{paragraph}{3}
\paragraph{}\etiquette[5.5.3.4]{XVII.par5.5.3.4}
若 $T$ \emph{为 $\leq n$ 次的}, \ie 第 $(n+1)$ 个“交叉效应”（cross-effect）为零 (\textsc{Dold-Puppe} \cite[4.18 \sisi{pg.232}{p.~232}]{sga4.exp.XVII.bib.DP})\sisi{}{,} 则对每个在 $i>m$ 时满足 $K^{i}=0$ 的 $K\in \Ob C^{\geq 0}(\Aa)$, 当 $i>nm$ 时有 $[T(K)]^{i}=0$ (\textsc{Dold-Puppe} \cite[4.2\sisi{.}{}3 \sisi{pg.233}{p.~233}]{sga4.exp.XVII.bib.DP}).

\paragraph{}\etiquette[5.5.3.5]{XVII.par5.5.3.5}
特别地, 若 $T$ 为从 $A$-模范畴到 $B$-模范畴的有限次函子, 并将平坦对象映为平坦对象, 则有界平坦复形的像仍是有界平坦复形.

\begin{enonce*}{引理 5.5.4}\etiquette[5.5.4]{XVII.lem5.5.4}
设 $\Aa$ 为阿贝尔范畴.
\begin{enumerate}
\renewcommand{\theenumi}{\roman{enumi}}
\renewcommand{\labelenumi}{(\theenumi)}
\item $K^{\geq 0}(\Aa)$ 中的拟同构容许左分式演算与右分式演算.

\item 显然的函子给出一个等价, 从 $K^{\geq 0}(\Aa)$ 中将拟同构变为可逆态射所得的范畴到范畴 $D^{\geq 0}(\Aa)$.
\end{enumerate}
\end{enonce*}

函子\pageoriginale $\tau_{\geq 0}:C(\Aa)\rightarrow
C^{\geq 0}(\Aa)$ (\hyperref[XVII.def1.1.13]{1.1.13}) 下降到商范畴后定义函子 $\tau_{\geq 0}:K(\Aa)\rightarrow K^{\geq
0}(\Aa)$, 它是 $K^{\geq
0}(\Aa)$ 到 $\sisi{\Aa}{K(\Aa)}$ 的包含函子的左伴随. 函子 $\tau_{\geq 0}$ 将拟同构映为拟同构, 并且伴随箭头 $X\rightarrow \tau_{\geq 0}(X)$ 为拟同构, 当且仅当在 $i<0$ 时 $H^{i}(X)=0$. 因而引理 \hyperref[XVII.lem5.5.4]{5.5.4} 来自 \textsc{Verdier} \cite{sga4.exp.XVII.bib.CD} 以及下述引理的 (i) 和 (ii); 其证明留给读者:

\begin{enonce*}{引理 5.5.4.1}\etiquette[5.5.4.1]{XVII.lem5.5.4.1}
设 $K$ 为范畴, $K'$ 为满子范畴, $S$ 为 $K$ 中的一组态射, 它等于 $K(S^{-1})$ 的同构集合的逆像 (\cf \cite[\sisi{chap 1}{\chap~1}]{sga4.exp.XVII.bib.GZ}); $S'$ 为它在 $K'$ 中的迹. 假设 $K'$ 到 $K$ 的包含函子 $i$ 有左伴随 $\tau$, 且 $\tau(S)\subset S'$.
\begin{enumerate}
\renewcommand{\theenumi}{\roman{enumi}}
\renewcommand{\labelenumi}{(\theenumi)}
\item 从 $K'({S'}^{-1})$ 到 $K(S^{-1})$ 的显然函子为全忠实函子, 并以 $\tau$ 为左伴随.

\item 若 $K(S^{-1})$ 容许左分式演算 (\resp 右分式演算), 则 $K'({S'}^{-1})$ 也具有相同性质.

若包含函子 $i$ 未必有左伴随, 且以 $\tau$ 表示与 $i$ 左伴随的\sisi{部分}{不处处}定义的函子, 则有:

\item 若 $K(S^{-1})$ 容许左分式演算 (\resp 右分式演算), 并且对每个终点 (\resp 起点) 在 $K'$ 中的 $u\in S$, $\tau(u)$ 都有定义且属于 $S'$, 则 $K'({S'}^{-1})$ 容许左分式演算 (\resp 右分式演算), 且函子 $K'({S'}^{-1})\rightarrow K(S^{-1})$ 全忠实.
\end{enumerate}
\end{enonce*}

\setcounter{subsubsection}{4}
\subsubsection{}\etiquette[5.5.5]{XVII.subsubsec5.5.5}
设\pageoriginale $T:\Aa\rightarrow \Bb$ 为阿贝尔范畴之间的函子. 如 \hyperref[XVII.def1.2.1]{1.2.1} 中那样, 由公式
\begin{align*}
LT(K) &= \quotprojlim_{s:K'\rightarrow K}T(K')\\
RT(K) &= \quotinjlim_{s:K\rightarrow K'}T(K'){\quoi},
\end{align*}
定义 $T$ 的导出函子
\begin{align*}
LT &: D^{\geq 0}(\Aa)\longrightarrow \Pro(D^{\geq
0}(\Bb))\\
RT &: D^{\geq 0}(\Aa)\longrightarrow \Ind(D^{\geq 0}(\Bb))
\end{align*}
其中 $s$ 遍历终点 (\resp 起点) 为 $K$ 的拟同构所成的滤过范畴.

若 $LT(K)$ $(RT(K))$ 可表, 则称 $T$ \emph{在 $K$ 处左 (\resp 右) 可导}. 若下述典范箭头为同构, 则称 $K$ \emph{左 (\resp 右) 适合于} $T$ (或滥用术语, 称其适合于 $LT$ (\resp $RT$)):
\[
LT(K)\isomto  TK\text{ (\resp $TK \isomto  RT(K)$)}{\quoi}.
\]

\begin{enonce*}{命题 5.5.6}\etiquette[5.5.6]{XVII.prop5.5.6}
设 $F$ 为从环 $A$ 上的模范畴到阿贝尔范畴 $\Bb$ 的函子. 假设 $\Bb$ 中的滤过归纳极限可表且正合, 而 $F$ 与滤过归纳极限可交换. 则每个平坦模的有界复形 $M^{\bigdot}\in \Ob K^{\geq
0}(A)$ 都适合于 $LF$.
\end{enonce*}

命题\pageoriginale 直接来自定义、下述将于 \hyperref[XVII.par5.5.6.4]{5.5.6.4} 证明的引理, 以及下述事实: 分量平坦的 $N^{\bigdot}\in \sisi{B}{C}^{b,\geq 0}(A)$ 所给出的拟同构 $u:N^{\bigdot}\rightarrow M^{\bigdot}$, 在终点为 $M^{\bigdot}$ 的拟同构范畴中构成一个共尾集合.

\begin{enonce*}{引理 5.5.6.1}\etiquette[5.5.6.1]{XVII.lem5.5.6.1}
设 $u:N^{\bigdot}\rightarrow M^{\bigdot}$ 为平坦复形之间的拟同构, 这些复形是 $C^{b,\geq 0}(A)$ 的对象. 则 $F(u)$ 为拟同构.
\end{enonce*}

\begin{enonce*}{引理 5.5.6.2}\etiquette[5.5.6.2]{XVII.lem5.5.6.2}
设 $\Aa$ 为加性范畴, $d$ 为 $\ZZ$ 中的上有界区间. 以 $C^{d}$ (\resp $C^{d,b}$, \resp $K^{d}$, \resp $K^{d,b}$) 表示在次数 $i\not\in d$ 处为零的复形 (\resp 且有界, \resp 模去同伦) 所成的范畴. 有一个等价
$$
\sisi{\sheafell }{\ell}:\Ind(C^{d,b}(\Aa)\sisi{}{)}\isomto
C^{d}(\Ind\Aa){\quoi};
$$
这一等价的逆诱导函子
$$
\sisi{\sheafell }{\ell}':K^{d}(\Ind\Aa)\longrightarrow
\Ind(K^{d,b}(\Aa)){\quoi}.
$$
\end{enonce*}

\begin{enumerate}
\renewcommand{\theenumi}{\Alph{enumi}}
\renewcommand{\labelenumi}{\theenumi)}
\item 证明 $\sisi{\sheafell }{\ell}$ 全忠实. 设 $K$ 为 $\Aa$ 中对象的有界复形, $L=\quotinjlim
L_{\alpha}$ 为 \iInd-复形. 有
\begin{align*}
\Hom(K,L)\overset{\dfn}{=} \varinjlim \Hom(K,L_{\alpha})  = & \varinjlim
Z^{0}(\Hom^{\bigdot} (K,L_{\alpha}))\simeq Z^{0}
\varinjlim_{\alpha} \Hom^{\bigdot} (K,L_{\alpha})\\
& \isomto
Z^{0} \Hom^{\bigdot} (K,\ell L)=\Hom(K,\ell L){\quoi},
\end{align*}
倒数第二个同构来自 $K$ 有限这一事实, 因而 $\varinjlim$ 与从双复形过渡到全复形可交换. 断言得证.

\item 我们\pageoriginale 将验证 $\sisi{\sheafell }{\ell}$ 本质满射的 (颇繁琐的) 工作留给读者: 从任意选定的次数开始, 从左向右构造映到一个上有界 \iInd-对象复形 $K$ 的“足够多”的有界复形; 同样, 留给读者展开如下论证:

\item 若 $\ell f:\ell \varinjlim K_{\alpha}\rightarrow \ell \varinjlim
L_{\beta}$ 零伦, $f=dH+Hd$, 则对每个 $\alpha$, $H$ 都提升为 $H^{i}_{\alpha}:K^{i}_{\alpha}\rightarrow
L^{i\sisi{}{-1}}_{\beta(\alpha)}$, 而当 $\beta$ 足够大时, 它给出 $f_{\alpha}:K_{\alpha}\rightarrow L_{\beta(\alpha)}$ 的零伦.
\end{enumerate}

\setcounter{subsubsection}{6}
\setcounter{paragraph}{2}
\paragraph{}\etiquette[5.5.6.3]{XVII.par5.5.6.3}
设 $A$ 为环, $d$ 为 $\ZZ$ 中的上有界区间, $P(A)$ 为有限投射 $A$-模的范畴, $\Pl(A)$ 为平坦 $A$-模的范畴. 以 $D^{\Tor\subset d}(A)$ 表示 $C^{d}(\Pl(A))$ 在 $D(A)$ 中的像, 以 $D^{\parf\subset d}(A)$ 表示 $C^{d,b}(P(A))$ 在 $D(A)$ 中的像. 由 \hyperref[XVII.par5.5.1.2]{5.5.1.2}, 有 $\Pl(A)\simeq \Ind
P(A)$. 另一方面, $K^{d,b}(P(A))\simeq D^{\parf\subset d}(A)$. 于是由 \hyperref[XVII.lem5.5.6.2]{5.5.6.2} 以及 $H^{i}$ 与滤过归纳极限可交换这一事实, 存在函子的交换图
$$
\xymatrix@C=2.25cm{
C^{d}(\text{Pl}(A)) \ar[d]& \Ind C^{d,b}(P(A)) \ar[d]\ar[l]_{\substack{\varinjlim\\\sim}}\\
K^{d}(\text{Pl}(A)) \ar[d]\ar[r]^{\ell'}& \Ind K^{d,b}(P(A)) \ar[d]^{\sisi{}{\simeq}}\\
D^{\Tor\subset d}(A) \ar[r]^{\ell''}& \Ind D^{\parf \subset d}(A).
}
$$

\setcounter{subsubsection}{6}
\paragraph{}\etiquette[5.5.6.4]{XVII.par5.5.6.4}
我们证明\pageoriginale \hyperref[XVII.lem5.5.6.1]{5.5.6.1}. 用 \hyperref[XVII.par5.5.6.3]{5.5.6.3} 的记号, 若 $d$ 为区间 $[0,k]$, 则 (由于 $F$ 与 $\varinjlim$ 可交换) 有交换图
$$
\xymatrix@C=2.5cm{
K^{d}(\text{Pl}(A))          \ar[r]^{\ell'} \ar[d]_{K^{\geq 0}(F)} & \Ind K^d(P(A))  \ar[d]_{\Ind(K^{\geq 0}(F))}             \\
K^{+}(\Bb) \ar[d]_{H^i}                        & \Ind K^{+}(\Bb)     \ar[d]_{\Ind(H^i)}       \\
\Bb                  & \Ind \Bb  \ar[l]_{\varinjlim}.
}
$$
若 $u$ 如 \hyperref[XVII.lem5.5.6.1]{5.5.6.1}, 则由 \hyperref[XVII.par5.5.6.3]{5.5.6.3}, $\ell'(u)$ 为同构, 因而 $H^iK^{\geq0}(F)(u)$ 也为同构, 如所欲证.

\subsubsection{}\etiquette[5.5.7]{XVII.subsubsec5.5.7}
设 $X$ 为概形 $S$ 上的概形, $k$ 为整数. 假设对称幂 $\Sym^k_S(X)$ 有定义. 若读者不愿作这一假设, 可在代数空间范畴中进行 (\textsc{Artin}\sisi{ [1]}{\nde{原文此处引用了一篇未指明的著作, 可能是第 \pageref{XVII.noteArtin1} 页和第 \pageref{XVII.noteArtin2} 页注释中所引的 Artin 的两篇著作之一.}}). 以\emph{交换}环层 $\Aa$ 给 $S$ 赋环. 从 \hyperref[XVII.subsubsec5.5.19]{5.5.19} 开始, 将假设 $\Aa$ 为\emph{挠的}. 滥用记号, 我们\emph{仍以} $\Aa$ 表示 $\Aa$ 在任一 $S$-概形上的逆像.

设 $ℱ$ 为 $X$ 上的模层; $ℱ$ 的 $k$ 份外张量积所成的层 $\mathop{\boxtimes}\limits^kℱ$ 是 $X^k/S$ 上的层, 对称群 $\Sf_k$ 在 $X^k/S$ 上作用时它具有等变性. 取不变量后,\pageoriginale 它定义 $\Sym^k_S(X)$ 上的一个层 $TS^k_{\ext}(ℱ)$, 称为 \emph{$ℱ$ 的外对称张量幂}（external symmetric tensor power）\sisi{}{.}

为了“逐点”论证, 采用如下描述的 $\Sym^k_S(X)$ 的点是有用的.

\paragraph{}\etiquette[5.5.7.1]{XVII.par5.5.7.1}
设 $s$ 为 $S$ 的一个几何点. $\Sym^k_S(X)$ 的一个取值于 $k(s)$ 的点 $x$ 定义 $\Sym^k_S(X)$ 的一个几何点. 以 $(x_1\ldots x_k)$ 表示 $(X/S)^k$ 的一个取值于 $k(s)$ 且像为 $x$ 的点, 设 $y_1\ldots y_\ell$ 为序列 $x_1\ldots x_k$ 中互异的元素, $n_i$ $(1\leq i\leq\ell)$ 为 $y_i$ 在该序列中的重数 $(\sum n_i=k)$.

用 \hyperref[XVII.par5.5.7.1]{5.5.7.1} 的记号, 并以 $\Sf$ 表示对称群, 有
$$\etiquette[5.5.7.2]{XVII.eq5.5.7.2}
TS^{k}_{\ext}(F)_x\simeq (\bigotimes^{\ell}_{1}
(\bigotimes^{n_i} F_{y_i}))^{\prod^{\ell}_{1}\Sf_{n_i}}{\quoi},\leqno(5.5.7.2)
$$
其中张量积在 $\Aa_S$ 上取. 特别地, \emph{当 $F$ 平坦时}, 由 \hyperref[XVII.par5.5.2.5]{5.5.2.5}, 有
$$\etiquette[5.5.7.3]{XVII.eq5.5.7.3}
TS^{k}_{\ext}(F)_x\isomfrom \bigotimes^{\ell}_{1}
\Gamma^{n_i}F_{y_i}{\quoi},\leqno(5.5.7.3)
$$
这里对 $\ell$ 归纳, 利用 $\Gamma^{n_i}F_{y_i}$ 的平坦性 (\hyperref[XVII.par5.5.2.4]{5.5.2.4}) 以及下述事实: 若 $G$ (\resp $H$) 为作用于 $A$-模 $M$\sisi{}{ }(\resp $N$) 的有限群, 且 $M$ 与 $N^H$ 平坦, 则 $M^G\otimes_A N^H\isomto(M\otimes_A N)^{G\times H}$.

\subsubsection{}\etiquette[5.5.8]{XVII.subsubsec5.5.8}
设 $F$ 为 $X$ 上的模, $P$ 为 $F$ 的一个平坦表示:
$$
P:P_1\longrightarrow P_0\longrightarrow F\longrightarrow0{\quoi}.
$$
如 \hyperref[XVII.par5.5.2.5]{5.5.2.5} 中那样定义\pageoriginale $P'_1$, $j_0$ 及 $j_1$; 将 $\Gamma^k_{\ext}(F)$ (即\emph{外 $\Gamma$ 函子}) 定义为一对箭头 $(TS^k_{\ext}(j_0),TS^k_{\ext}(j_1))$ 的余核. 用 \hyperref[XVII.par5.5.7.1]{5.5.7.1} 的记号, 由 \eqhyperref[XVII.eq5.5.7.3]{5.5.7.3} 及 \hyperref[XVII.par5.5.2.5]{5.5.2.5} 得到同构
$$\etiquette[5.5.8.1]{XVII.eq5.5.8.1}
\Gamma^k_{\ext}(F)_x\simeq\bigotimes^{\ell}_{1}
\Gamma^{n_i}(F_{y_i}){\quoi}.\leqno(5.5.8.1)
$$

这一公式表明, $\Gamma^k_{\ext}(F)$ 先验地函子依赖于 $P$, 但在相差唯一同构的意义下只依赖于 $F$.

\subsubsection{}\etiquette[5.5.9]{XVII.subsubsec5.5.9}
存在函子性态射
$$
\alpha:\Gamma^k_{\ext}(F)\longrightarrow TS^k_{\ext}(F){\quoi}.
$$
对 $(\Gamma,\alpha)$ 由下列两项性质刻画:
\begin{enumerate}
\renewcommand{\theenumi}{\roman{enumi}}
\renewcommand{\labelenumi}{(\theenumi)}
\item 当 $F$ 平坦时, $\alpha$ 为同构.

\item 用 \hyperref[XVII.par5.5.2.5]{5.5.2.5} 的记号, 对 $F$ 的每个表示
$$
P_1\xrightarrow{d}P_0\xrightarrow{\epsilon}F\longrightarrow0
$$
序列
$$
\Gamma^k_{\ext}(P'_1)\rightrightarrows
\Gamma^k_{\ext}(P_0)\rightarrow\Gamma^k_{\ext}(F)\rightarrow0
$$
都正合.

此外:

\item 函子 $TS^k$ 与 $\Gamma^k$ 都与滤过归纳极限可交换 (由 \sisi{\hyperref[XVII.eq5.5.7.2]{5.5.7.2}}{\eqhyperref[XVII.eq5.5.7.2]{5.5.7.2}}, \sisi{\hyperref[XVII.eq5.5.8.1]{5.5.8.1}}{\eqhyperref[XVII.eq5.5.8.1]{5.5.8.1}}, \hyperref[XVII.par5.5.2.4]{5.5.2.4}).

\item 函子 $\Gamma^k$ 与底环的每个扩张 $\Aa\rightarrow\Aa'$ 可交换.

\item 函子\pageoriginale $\Gamma^k$ 与每个基变换 $S'\rightarrow S$ 可交换: 对每个笛卡儿方块
$$
\xymatrix@C=2.5cm@R=1.5cm{
X' \ar[r]^{g'}\ar[d]^{f'}& X \ar[d]^{f}\\
S' \ar[r]^{g}& S,
}
$$
有
$$
\Sym^k_g(g')^{\ast}\Gamma^k_{\ext}\simeq
\Gamma^k_{\ext}{g'}^{\ast}{\quoi}.
$$
(\sisi{\hyperref[XVII.eq5.5.8.1]{5.5.8.1}}{\eqhyperref[XVII.eq5.5.8.1]{5.5.8.1}} 及 \hyperref[XVII.par5.5.2.7]{5.5.2.7}).

\item 函子 $\Gamma^k$ 与沿每个\emph{根}态射（radicial morphism）$i:X'\rightarrow X$ (例如闭浸入) 的逆像可交换: 有
$$
\Sym^k_S(i)^{\ast}\Gamma^k_{\ext}\simeq
\Gamma^k_{\ext}i^{\ast}
$$
\eqhyperref[XVII.eq5.5.8.1]{5.5.8.1}.

\item 对\emph{浸入} $j:X'\rightarrow X$, 有
$$
\Sym^k_S(j)_{!}\Gamma^k_{\ext}\simeq\Gamma^k_{\ext}j_{!}
$$
\eqhyperref[XVII.eq5.5.8.1]{5.5.8.1}.
\end{enumerate}

\subsubsection{}\etiquette[5.5.10]{XVII.subsubsec5.5.10}

设 $k_1,k_2$ 及 $k=k_1+k_2$ 为三个整数 $\geq0$. 则有有限态射 (\quad):\nde{我未找到一般情形下 $\vee_{k_1k_2}$ 有限性的参考文献.}
$$
\vee_{k_1k_2}:\Sym^{k_1}_S(X)\times_S\Sym^{k_2}_S(X)
\longrightarrow\Sym^k_S(X){\quoi}.
$$
若 $F_i$ $(i=1,2)$ 为 $\Sym_S^{k_i}(X)$ 上的模, 以 $F_1\vee F_2$ 表示 $\Sym_S^k(X)$ 上的层
$$
F_1\vee F_2=\vee_{k_1k_2\ast}(F_1\boxtimes F_2){\quoi}.
$$

\subsubsection{}\etiquette[5.5.11]{XVII.subsubsec5.5.11}
设\pageoriginale $X$ 上有模的正合列
\[
\Sigma:0\longrightarrow F'\longrightarrow F\longrightarrow
F''\longrightarrow0{\quoi}.
\]
用 \hyperref[XVII.par5.5.7.1]{5.5.7.1} 的记号, 由 \eqhyperref[XVII.eq5.5.8.1]{5.5.8.1} 及 \hyperref[XVII.par5.5.2.6]{5.5.2.6}, $\Gamma^k_{\ext}(F)_x$ 上有一个滤过, 并有态射
\[
\sum_{\sum^{\ell}_{1}i_\alpha=i}
\bigotimes^{\ell}_{1}\Gamma^{i_\alpha}(F'_{y_\alpha})\otimes
\bigotimes^{\ell}_{1}\Gamma^{n_i-i_\alpha}(F''_{y_\alpha})
\longrightarrow\Gr^i(\Gamma^k_{\ext}(F)_x){\quoi}.
\]
第一项恰为 $(\Gamma^i_{\ext}(F')\vee\Gamma^{k-i}_{\ext}(F''))_x$. 容易验证, 这些逐点构造的数据来自 $\Gamma^k_{\ext}(F)$ 上的 (唯一) $k+1$ 步滤过, 以及 (唯一的) 态射
$$\etiquette[5.5.11.1]{XVII.eq5.5.11.1}
\Gamma^i_{\ext}(F')\vee\Gamma^{k-i}_{\ext}(F'')
\longrightarrow\Gr^i(\Gamma^k_{\ext}(F)){\quoi}.\leqno(5.5.11.1)
$$

若正合列 $\Sigma$ 分裂, 更准确地有 (由 \sisi{\hyperref[XVII.eq5.5.2.6.2]{5.5.2.6.2}}{\eqhyperref[XVII.eq5.5.2.6.1]{5.5.2.6.1}} 逐点验证):
$$\etiquette[5.5.11.2]{XVII.eq5.5.11.2}
\Gamma^k_{\ext}(F'\oplus F'')\simeq
\sisi{\sum}{\bigoplus}_{k_1+k_2=k}\Gamma^{k_1}_{\ext}(F')\vee
\Gamma^{k_2}_{\ext}(F''){\quoi}.\leqno(5.5.11.2)
$$

若正合列 $\Sigma$ 逐点分裂, 或者 $F''$ 平坦, 则态射 \eqhyperref[XVII.eq5.5.11.1]{5.5.11.1} 为同构 (\cf \hyperref[XVII.par5.5.1.4]{5.5.1.4}).

\subsubsection{}\etiquette[5.5.12]{XVII.subsubsec5.5.12}
设 $u:X\rightarrow Y$ 为 $S$-概形态射. 假设对称幂 $\Sym^k_S(X)$ 与 $\Sym^k_S(Y)$ 存在. 设 $F$ 为 $X$ 上的模, $G$ 为 $Y$ 上的模, 并设
$v\in\Hom(u^{\ast}G,F)\underset{\sisi{n}{}}{\simeq}\Hom(G,u_{\ast}F)$
为从 $F$ 到 $G$ 的 $u$-同态 \hyperlink{sga4.exp.XVII.subsection.0.2}{(0.2)}. 若 $\prod^n u$ 为 $u$ 的第 $n$ 笛卡儿幂, 则同态 $v$ 诱导一个从 $\mathop{\boxtimes}\limits^nF$ 到 $\mathop{\boxtimes}\limits^nG$ 的 $\prod^n u$-同态 $\mathop{\boxtimes}\limits^n v$. 取不变量后,\pageoriginale 后者定义
\[
TS^n(v):TS^n_{\ext}(F)\xrightarrow[\Sym_S^n(u)]{}TS^n_{\ext}(G){\quoi}.
\]

由定义 \hyperref[XVII.subsubsec5.5.8]{5.5.8}, 存在唯一的 $\Sym^n(u)$-同态 $\Gamma^n_{\ext}(v)$, 关于 $(F,G,v)$ 具有函子性,
$$\etiquette[5.5.12.0]{XVII.eq5.5.12.0}
\Gamma^n_{\ext}(v):\Gamma^n_{\ext}(F)
\xrightarrow[\Sym^n_S(u)]{}\Gamma^n_{\ext}(G)\leqno(5.5.12.0)
$$
使下图交换 (关于箭头的不寻常方向, \cf \hyperlink{sga4.exp.XVII.subsection.0.2}{(0.2)}):
$$
\xymatrix@C=3.5cm@R=2cm{
\Gamma^n_{\ext}(F) \ar[r]^{\Gamma^n_{\ext}(v)}_{\text{在 $\Sym^n_S(u)$ 上}}& \Gamma^n_{\ext}(G)\\
TS^n_{\ext}(F) \ar[r]^{TS^n_{\ext}(v)}_{\text{在 $\Sym^n_S(u)$ 上}}\ar[u]^{\alpha}_{\text{在 $\Id$ 上}} & TS^n_{\ext}(G)\ar[u]^{\alpha}_{\text{在 $\Id$ 上}}.
}
$$

我们将主要把这条箭头用作“对称 Künneth 箭头”: 对 $G=u_{\ast}F$ 以及典范箭头 $v$, 箭头 $\Gamma^n_{\ext}(v)$ 等同于模的同态\sisi{s}{}
$$\etiquette[5.5.12.1]{XVII.eq5.5.12.1}
k^n:\Gamma^n_{\ext}(u_{\ast}F)\longrightarrow
\Sym^n(u)_{\ast}\Gamma^n_{\ext}(F){\quoi}.\leqno(5.5.12.1)
$$

\subsubsection{}\etiquette[5.5.13]{XVII.subsubsec5.5.13}
考虑条件:

\paragraph{}\etiquette[5.5.13.1]{XVII.par5.5.13.1}
$K\in\Ob K^{b,\geq0}(X,\Aa)$, 并且对 $X$ 的每个几何点 $x$, 复形 $K_x$ 同伦等价于一个平坦模复形 $L_x\in\Ob C^{b,\geq0}(\Aa_x)$.

\begin{enonce*}{引理 5.5.13.2}\etiquette[5.5.13.2]{XVII.lem5.5.13.2}
$D^b(X,\Aa)$ 中由 \tTor-维数 $\leq0$ \emph{\hyperref[XVII.def4.1.9]{4.1.9}} 的复形组成的子范畴\pageoriginale $D^{b,\tor\leq0}(X,\Aa)$, 等价于将满足\emph{(\hyperref[XVII.par5.5.13.1]{5.5.13.1})} 的复形 $K\in\Ob K^{b,\geq0}(X,\Aa)$ 的范畴中的拟同构变为可逆态射所得的范畴. 这一构造可以同样用左分式演算或右分式演算进行.
\end{enonce*}

设 $P$ 为 $K^{-}(X,\Aa)$ 中由逐点同伦等价于平坦模复形的复形组成的子范畴. 已知 (\textsc{Verdier} \sisi{\cite{sga4.exp.XVII.bib.CD}}{\cite[I 2.4.2]{sga4.exp.XVII.bib.CD}}), $D^{-}(X,\Aa)$ 是从 $P$ 通过左分式演算或右分式演算将拟同构变为可逆态射而得到的. 将 \hyperref[XVII.lem5.5.4.1]{5.5.4.1} 应用于 $K^{b,\geq0}(X,\Aa)\cap P$ 到 $P$ 的包含函子. 事实上, 若 $K$ 的 \tTor-维数为 $\leq0$, 并逐点同伦等价于平坦模复形 $L_x$, 则 $\tau_{\geq0}(K)$ 逐点同伦等价于 $\tau_{\geq0}(L_x)$, 而后者仍为平坦模复形, 因而 $\tau_{\geq0}(K)$ 满足 \hyperref[XVII.par5.5.13.1]{5.5.13.1}.

\subsubsection{}\etiquette[5.5.14]{XVII.subsubsec5.5.14}
由 \hyperref[XVII.lem5.5.6.1]{5.5.6.1}, \sisi{\hyperref[XVII.eq5.5.7.2]{5.5.7.2}}{\eqhyperref[XVII.eq5.5.7.2]{5.5.7.2}} 或 \sisi{\hyperref[XVII.eq5.5.8.1]{5.5.8.1}}{\eqhyperref[XVII.eq5.5.8.1]{5.5.8.1}} 以及 \hyperref[XVII.par5.5.2.4]{5.5.2.4}, 满足 (\hyperref[XVII.par5.5.13.1]{5.5.13.1}) 的复形之间的每个拟同构 $v:K'\rightarrow K$ 都诱导拟同构
\begin{align*}
TS^k(v)&:TS^k_{\ext}(K')\xrightarrow{\approx}TS^k_{\ext}(K)\\
\Gamma^k(v)&:\Gamma^k_{\ext}(K')\xrightarrow{\approx}\Gamma^k_{\ext}(K){\quoi}.
\end{align*}
因此函子 $\Gamma^k$ 与 $TS^k$ 在 $D^b(X,\Aa)$ 的子范畴 $D^{b,\tor\leq0}(X,\Aa)$ 上可导 (\hyperref[XVII.subsubsec5.5.5]{5.5.5}), 而满足 \hyperref[XVII.par5.5.13.1]{5.5.13.1} 的复形左适合于这些函子. 由 \hyperref[XVII.par5.5.2.4]{5.5.2.4},\pageoriginale $\Gamma^k$ 与 $TS^k$ 诱导函子
\begin{align*}
LTS^k &:D^{b,\tor\leq0}(X,\Aa)\longrightarrow
D^{b,\tor\leq0}(\Sym^k_S(X),\Aa)\tag{5.5.14.1}\\
L\Gamma^k &:D^{b,\tor\leq0}(X,\Aa)\longrightarrow
D^{b,\tor\leq0}(\Sym^k_S(X),\Aa){\quoi}.\tag{5.5.14.2}
\end{align*}\etiquette[5.5.14.1]{XVII.eq5.5.14.1}
\etiquette[5.5.14.2]{XVII.eq5.5.14.2}% 删除此注释会引入不需要的空白.
此外, 由 \hyperref[XVII.subsubsec5.5.9]{5.5.9}~(i), 典范态射 $\alpha$ (\hyperref[XVII.subsubsec5.5.9]{5.5.9}) 在函子 \eqhyperref[XVII.eq5.5.14.1]{5.5.14.1} 与 \eqhyperref[XVII.eq5.5.14.2]{5.5.14.2} 之间诱导同构:
$$\etiquette[5.5.14.3]{XVII.eq5.5.14.3}
\alpha:L\Gamma^k\isomto LTS^k{\quoi}.\leqno(5.5.14.3)
$$

\setcounter{paragraph}{3}
\paragraph{}\etiquette[5.5.14.4]{XVII.par5.5.14.4}
\hyperref[XVII.subsubsec5.5.9]{5.5.9}~(iv), (v), (vi) 及 (vii) 中的同构过渡到导出范畴后定义同构:
\begin{itemize}
\item[${\rm(iv')}$] 对 $\Aa$-代数 $\Aa'$, 有
$$
L\Gamma^k_{\ext_{\Aa}}(K)\overset{\LL}{\otimes}\Aa'\simeq
L\Gamma^k_{\ext_{\Aa'}}(K\overset{\LL}{\underset{\Aa}{\otimes}}\Aa'){\quoi}.
$$

\item[${\rm(v')}$] 对每个基变换 $g:S'\rightarrow S$, 有
$$
\Sym^k_g(g')^{\ast}L\Gamma^k_{\ext}K\simeq
L\Gamma^k_{\ext}({g'}^{\ast}K){\quoi}.
$$

\item[${\rm(vi')}$] 对每个\emph{根}态射 $i:X'\rightarrow X$, 有
$$
\Sym^k_S(i)^{\ast}L\Gamma^k_{\ext}K\simeq
L\Gamma^k_{\ext}i^{\ast}K{\quoi}.
$$

\item[${\rm(vii')}$] 对\emph{浸入} $j:X'\rightarrow X$, 有
$$
\Sym^k_S(j)_{!}L\Gamma^k_{\ext}\simeq
L\Gamma^k_{\ext}j_{!}{\quoi}.
$$
\end{itemize}

\subsubsection{}\etiquette[5.5.15]{XVII.subsubsec5.5.15}
设 $C^{\geq0}(X,\Aa)$ 中有平坦复形的正合列,
$$
0\longrightarrow K'\longrightarrow K\longrightarrow K''\longrightarrow0{\quoi},
$$
并应用函子 $\Gamma^k_{\ext}$ (\hyperref[XVII.subsubsec5.5.3]{5.5.3}). 可见 (如 \hyperref[XVII.subsubsec5.5.11]{5.5.11} 中一样), $\Sym^k_S(X)$ 上的半单纯模 $\Gamma^k_{\ext}(N^{-1}K)$ 有一个滤过, 其\pageoriginale 逐次商为半单纯模 $\Gamma^{k'}_{\ext}(N^{-1}K')\vee
\Gamma^{k''}_{\ext}(N^{-1}K'')$ $(k'+k''=k)$.

过渡到导出范畴可见, $L\Gamma^k_{\ext}(K)$ 由一个滤过复形表示, 其逐次商表示导出范畴的对象 $L\Gamma^i_{\ext}(K'\sisi{}{)}\overset{\LL}{\vee}
L\Gamma^{k-i}_{\ext}(K'')$ (应用 \hyperref[XVII.par5.5.2.4]{5.5.2.4} 及 \textsc{Eilenberg--Zilber--Cartier} 定理, 即双半单纯对象与其对角对象在 $\mathcal K$ 中同构: \textsc{Dold\sisi{ }{-}Puppe} \cite[2.9 \sisi{pg.213}{p.~213}]{sga4.exp.XVII.bib.DP}).

在导出范畴中, 从满足 \hyperref[XVII.par5.5.13.1]{5.5.13.1} 的复形的一个正合列出发, 已经可以进行这一构造.

${}^{\ast}$ 用更内在的术语说, 设 $\Delta$ 为 $D^{b,\tor\leq0}(X,\Aa)$ 的对象的一个“\emph{真三角}”, \ie 为“滤过导出范畴”的一个对象, 在 $i\neq0,1$ 时满足 $\Gr^i(\Delta)=0$, 且 $\Gr^i(\Delta)\in\Ob D^{b,\tor\leq0}(X,\Aa)$ (\textsc{Deligne}, 见 \textsc{Verdier} \cite{sga4.exp.XVII.bib.V2} 的附录, 或 \refvb[VI B nᵒ]{Vbis.subsubsec2.5.1}, \refnoexpose{Vbis.subsubsec2.5.2}).

前述构造于是给出 $\Sym^k_S(X)$ 的滤过导出范畴的对象 $L\Gamma^k_{\ext}(\Delta)$、同构
\[
\Gr^i(L\Gamma^k_{\ext}(\Delta))\simeq
L\Gamma^i_{\ext}(\Gr^1(\Delta))\overset{\LL}{\vee}
L\Gamma^{k-i}_{\ext}(\Gr^0(\Delta)\sisi{}{)}
\]
以及一个同构 (以 $\lvert\quad\rvert$ 表示滤过复形的底层复形)
\[
\lvert L\Gamma^k_{\ext}(\Delta)\rvert\simeq
L\Gamma^k_{\ext}(\lvert\Delta\rvert){\quoi}.{}_{\ast}
\]

\begin{enonce*}{引理 5.5.16}\etiquette[5.5.16]{XVII.lem5.5.16}
设 $N$ 为整数. 范畴 $D^{b,\tor\leq0}(X,\Aa)$ 可由满足下列条件 {\rm(i)} 和 {\rm(ii)} 的复形 $K\in\Ob K^{b,\geq0}(X,\Aa)$ 的范畴, 通过将拟同构变为可逆态射而得到:
\begin{enumerate}
\renewcommand{\theenumi}{\roman{enumi}}
\renewcommand{\labelenumi}{(\theenumi)}
\item 条件\pageoriginale (\hyperref[XVII.par5.5.13.1]{5.5.13.1}).

\item 复形 $K$ 适合于每个关于拓扑斯态射 $f:\widetilde X_{\text{\rm et}}\rightarrow T$ 的上同调维数 $\leq N$ 的正像函子.
\end{enumerate}
\end{enonce*}

右分式演算的存在性来自 \hyperref[XVII.subsubsec5.5.13]{5.5.13} 以及以下事实: 利用截断的典范松软消解（canonical flasque resolution）(\hyperref[XVII.subsubsec4.2.9]{4.2.9}), 每个满足 (i) 的复形 $K$ 都映到一个满足 (i) 和 (ii) 的复形.

\setcounter{subsubsection}{16}

\subsubsection{}\etiquette[5.5.17]{XVII.subsubsec5.5.17}
设 $u:X\rightarrow Y$ 为 $S$-概形态射. 假设 $\Sym^n_S(X)$ 与 $\Sym^n_S(Y)$ 存在, 存在整数 $N$ 使函子 $u_{\ast}$ 在 $\Aa$-模范畴上的维数为 $\leq N$, 且函子 $Ru_{\ast}$ 将 \tTor-维数 $\leq0$ 的有界复形映为 \tTor-维数 $\leq0$ 的复形. 若 $u$ 固有, $\Aa$ 为挠的, 且 $Y$ 凝聚 (\ie 拟紧且拟分离), 则后两项条件自动满足 (\hyperref[XVII.cor5.2.8.1]{5.2.8.1}, \hyperref[XVII.thm5.2.10]{5.2.10}).\nde{$Y$ 的凝聚性可以换成较弱的条件, 即 $u$ 的纤维维数有界.}

设 $K\in\Ob K^{b,\geq0}(X,\Aa)$ 满足 \hyperref[XVII.lem5.5.16]{5.5.16} 的条件 (i) 与 (ii), 其中 $N$ 如上. 由 \hyperref[XVII.subsubsec5.5.14]{5.5.14}, 箭头
$$
\beta:L\Gamma^n_{\ext}(K)\isomto\Gamma^n_{\ext}(K)
$$
为拟同构. 同样, 箭头
$$
\gamma:u_{\ast}K\isomto Ru_{\ast}K
$$
为拟同构.

将\pageoriginale 从 $L\Gamma^n_{\ext}Ru_{\ast}(K)$ 到 $R\Sym^n(u)_{\ast}L\Gamma^n_{\ext}(K)$ 的\emph{对称 Künneth 箭头}定义为复合态射
\begin{multline*}
k^{n}: L\Gamma^{n}_{\ext}Ru_{\ast}K\xleftarrow[\gamma]{\sim}
L\Gamma^{n}_{\ext}u_{\ast}K\rightarrow
\Gamma^{n}_{\ext}u_{\ast}K\xrightarrow[\protect\hyperlink{sga4.semantic.XVII.eq5.5.12.1}{5.5.12.1}]{}\Sym^{n}_{S}(u)_{\ast}\Gamma^{n}_{\ext}(K)\\
\longrightarrow R
\Sym^{n}_{S}(u)_{\ast}\Gamma^{n}_{\ext}(K)\xleftarrow[\beta]{\sim} R \Sym^{n}_{S}(u)_{\ast}L\Gamma^{n}_{\ext}(K){\quoi}.\tag{5.5.17.1}
\end{multline*}\etiquette[5.5.17.1]{XVII.eq5.5.17.1}% 本段之后的竖直间距过大.

由 \hyperref[XVII.lem5.5.16]{5.5.16}, 态射 \eqhyperref[XVII.eq5.5.17.1]{5.5.17.1} 定义从 $D^{b,\tor\leq0}(X,\Aa)$ 到 $D^{+}(\Sym^n_S(Y),\Aa)$ 的函子之间的态射:
$$\etiquette[5.5.17.2]{XVII.eq5.5.17.2}
k^n:L\Gamma^n_{\ext}Ru_{\ast}\longrightarrow
R\Sym^n_S(u)_{\ast}L\Gamma^n_{\ext}{\quoi},\leqno(5.5.17.2)
$$
也称为\emph{对称 Künneth 态射}.

这一态射由下述命题 \hyperref[XVII.prop5.5.18]{5.5.18} 刻画.

\setcounter{paragraph}{2}
\paragraph{}\etiquette[5.5.17.3]{XVII.par5.5.17.3}
设 $S$, $X$, $Y$ 及 $u$ 如上, $K\in\Ob C^{b,\geq0}(X,\Aa)$ 为 \tTor-维数 $\leq0$ 的复形, $K_1\in\Ob C^{b,\geq0}(Y,\Aa)$ 且 $K_2\in\Ob C^b(\Sym^n_S(X),\Aa)$. 设 $v$ 为从 $K$ 到 $K_1$ 的 $u$-同态 \hyperlink{sga4.exp.XVII.subsection.0.2}{(0.2)}, $t$ 为从 $\Gamma^n_{\ext}(K)$ 到 $K_2$ 的同态 (是通常意义的同态, 以下所有同态也如此). 有态射
\begin{enumerate}
\renewcommand{\theenumi}{\alph{enumi}}
\renewcommand{\labelenumi}{\theenumi)}
\itemsep=8pt
\item $\Gamma^n_{\ext}(v):\Gamma^n_{\ext}(K_1)\rightarrow
\Sym^n_S(u)_{\ast}\Gamma^n_{\ext}(K)$\qquad
\eqhyperref[XVII.eq5.5.12.0]{5.5.12.0}

\item
$\Sym^n_S(u)_{\ast}(t):\Sym^n_S(u)_{\ast}\Gamma^n_{\ext}(K)
\rightarrow\Sym^n_S(u)_{\ast}K_2$,
\end{enumerate}
它们的复合
$$
\varphi':\Gamma^n_{\ext}(K_1)\longrightarrow
\Sym^n_S(u)_{\ast}K_2{\quoi},
$$
通过复合定义
$$
\varphi:L\Gamma^n_{\ext}(K_1)\longrightarrow
R\Sym^n(u)_{\ast}K_2{\quoi}.
$$
此外,\pageoriginale $v$ 与 $t$ 诱导
$$
\begin{cases}
v':K_1\longrightarrow Ru_{\ast}K\\
t':L\Gamma^n_{\ext}K\longrightarrow K_2{\quoi}.
\end{cases}
$$

于是:

\begin{enonce*}{命题 5.5.18}\etiquette[5.5.18]{XVII.prop5.5.18}
在 \hyperref[XVII.par5.5.17.3]{5.5.17.3} 的假设下, 下图在 $D^{+}(\Sym^n_S(Y),\Aa)$ 中交换:
\[
\xymatrix@C=4cm@R=1.75cm{
  L\Gamma^n_{\ext}Ru_{\ast}K \ar[r]^{k^n}_{\text{\textsuperscript{(对称 Künneth)}}}& R\Sym^n_S(u)_{\ast}L
\Gamma^n_{\ext}K \ar[d]_{R
  \Sym^n_S(u)_{\ast}(t')}\\
L\Gamma^n_{\ext}K_1 \ar[r]^{\varphi}\ar[u]^{L\Gamma^n_{\ext}(v')} & R\Sym^n_S(u)_{\ast}K_2.
}
\]
\end{enonce*}

证明与 \hyperref[XVII.par5.4.1.5]{5.4.1.5} 的证明完全类似, 留给读者.

\setcounter{subsubsection}{18}
\subsubsection{从现在起假设 $\Aa$ 为挠的}\etiquette[5.5.19]{XVII.subsubsec5.5.19}

设 $u:X\rightarrow Y$ 为可紧化的 $S$-概形态射:
\[
\xymatrix@C=3cm@R=1.75cm{
   & \bar{X} \ar[rd]^{\bar{u}}&\\
X \ar@{^{(}->}[ru]^{j} \ar[rr]^{u} \ar[rd]&         & Y \ar[ld]\\
   &   S     & .
}
\]
假设\pageoriginale $\Sym^n_S(X)$, $\Sym^n_S(\bar X)$ 及 $\Sym^n_S(Y)$ 存在. 对 $X$ 上的每个模 $F$, 都有
$$\etiquette[5.5.19.1]{XVII.eq5.5.19.1}
\Gamma^n_{\ext}(j_{!}F)\simeq
\Sym^n_S(j)_{!}\Gamma^n_{\ext}(F){\quoi},\leqno(5.5.19.1)
$$
这一同构立即过渡到导出范畴 \sisi{}{(\hyperref[XVII.par5.5.14.4]{5.5.14.4}~(vii'))}. 将从 $D^{b,\tor\leq0}(X,\Aa)$ \linebreak 到 $D^{b,\tor\leq0}(\Sym^n_S(Y),\Aa)$\nde{一般情况下如何证明 $\Sym^n_S(u)$ 为有限型? 这一条件是 $R\Sym^n_S(u)_!$ 有定义的必要条件.} 的函子之间的\emph{对称 Künneth 态射}
$$\etiquette[5.5.19.2]{XVII.eq5.5.19.2}
k^n_u\quad\mbox{\rm 或}\quad
k^n:L\Gamma^n_{\ext}Ru_{!}\longrightarrow
R\Sym^n_S(u)_{!}L\Gamma^n_{\ext}\leqno(5.5.19.2)
$$
定义为复合态射
\begin{align*}
L\Gamma^n_{\ext}Ru_{!}\stackrel{\dfn}{\simeq}L\Gamma^n_{\ext}R\bar u_{\ast}j_{!}
&\xrightarrow{\eqhyperref[XVII.eq5.5.17.2]{5.5.17.2}}
R\Sym^n_S(\bar u)_{\ast}L\Gamma^n_{\ext}j_!\\
&\xrightarrow[\sisi{}{\text{(\hyperref[XVII.par5.5.14.4]{5.5.14.4} (vii'))}}]{\sisi{\substack{\eqhyperref[XVII.eq5.5.19.1]{5.5.19.1}\\\sim}}{\sim}}
R\Sym^n_S(\bar u)_*\Sym^n_S(j)_{!}L\Gamma^n_{\ext}\\
&\xrightarrow{\sisi{\text{(5.5.13.5 (vii'))}}{}}
R\Sym^n_S(u)_{!}L\Gamma^n_{\ext}.
\end{align*}

\setcounter{subsubsection}{19}
\subsubsection{}\etiquette[5.5.20]{XVII.subsubsec5.5.20}
以下验证留给读者, 利用 \hyperref[XVII.prop5.5.18]{5.5.18}, 原则上它们都是容易的.

\paragraph{}\etiquette[5.5.20.1]{XVII.par5.5.20.1}
(内在性). 态射 \eqhyperref[XVII.eq5.5.19.2]{5.5.19.2} 与 $u$ 的紧化 $\bar u$ 的选取无关.

\paragraph{}\etiquette[5.5.20.2]{XVII.par5.5.20.2}
(传递性). 设有 $S$-概形的图
$$
X\xrightarrow{u}Y\xrightarrow{v}Z{\quoi},
$$
其中\pageoriginale $u$ 与 $v$ 拟射影; 假设 $\Sym^n_S(Z)$ 存在 (从而 $\Sym^n_S(Y)$ 与 $\Sym^n_S(X)$ 也存在 (\SGA 1 V 1.7, \EGA II 4.5.4)). 则下列函子图交换:
$$
\xymatrix@C=1.5cm@R=1.75cm{
L\Gamma^n_{\ext}R(vu)_{!} \ar@{=}[d] \ar[rr]^{k^n_{vu}}& & R\Sym^n_S(vu)_{!}L\Gamma^n_{\ext}\ar@{=}[d]\\
L\Gamma^n_{\ext}Rv_{!}Ru_{!} \ar[r]^{k^n_v}& R\Sym^n_S(v)_!
L\Gamma^n_{\ext}Ru_{!} \ar[r]^{k^n_u}& R\Sym^n_S(v)_{!}R\Sym^n_S(u)_{!}L\Gamma^n_{\ext}.
}
$$

\paragraph{}\etiquette[5.5.20.3]{XVII.par5.5.20.3}
对称 Künneth 态射与 \hyperref[XVII.subsubsec5.5.9]{5.5.9}~(iv), (v) 及 (vi) 类型的底或空间变换相容. 更准确地, 与态射 \sisi{5.5.13.5}{\hyperref[XVII.par5.5.14.4]{5.5.14.4}} ${\rm(iv')}$, ${\rm(v')}$ 及 ${\rm(vi')}$ 有如下相容性.
\begin{itemize}
\item[${\rm(iv'')}$] 对每个 $\Aa$-代数 $\Aa'$ 以及每个拟射影 $S$-态射 $u:X\rightarrow Y$, 图
\[\def\objectstyle{\scriptstyle}
\xymatrix@C=.8cm@R=1.75cm{
L\Gamma^n_{\ext}Ru_{!}(K)\overset{\LL}{\otimes}_{\Aa}\Aa'
\ar@{-}[r]^{\sim} \ar[d]^{k^n_u}&
L\Gamma^n_{\ext}(Ru_{!}(K)\overset{\LL}{\otimes}_{\Aa}\Aa')
\ar@{->}[r]^{\sim}&
L\Gamma^n_{\ext}Ru_{!}(K\overset{\LL}{\otimes}_{\Aa}\Aa') \ar[d]^{k^n_u}\\
R\Sym^n_S(u)_{!}L\Gamma^n_{\ext}(K)\overset{\LL}{\otimes}_{\Aa}\Aa'
\ar@{->}[r]^{\sim}&
R\Sym^n_S(u)_{!}(L\Gamma^n_{\ext}(K)\overset{\LL}{\otimes}_{\Aa}\Aa')
\ar@{-}[r]^{\sim}&
R\Sym^n_S(u)_{!}L\Gamma^n_{\ext}(K\overset{\LL}{\otimes}_{\Aa}\Aa')
}
\]
交换.

\item[${\rm(v'')}$] 对\pageoriginale 每个基变换 $g:S'\rightarrow S$, 它给出一个\sisi{笛卡儿}{方块均为笛卡儿方块的}图
$$
\xymatrix{
X' \ar[rrr]^{g'}\ar[rd]^{u'}\ar[dd]&   & & X \ar[rd]^{u}\ar[dd]|\hole &   \\
   & Y' \ar[ld]\ar[rrr]^{g''}&  & & Y \ar[ld]\\
S' \ar[rrr]_{g}&  &  & S &   ,
}
$$
完全类似于 ${\rm(iv'')}$ 的函子图交换:
$$\def\objectstyle{\scriptstyle}
\xymatrix@R=2.25cm{
\Sym^n_g(g'')^{\ast}L\Gamma^n_{\ext}Ru_{!} \ar@{-}[r]^{\sim} \ar[d]^{k^n_u}&
L\Gamma^n_{\ext}g''^{\ast}Ru_{!} \ar[r]^{\sim}& L\Gamma^n_{\ext}Ru'_{!}
g'^{\ast}\ar[d]^{k^n_{u'}} \\
\Sym^n_g(g'')^{\ast}R\Sym^n_S(u)_{!}L\Gamma^n_{\ext} \ar[r]^{\sim}& R
\Sym^n_S(u')_{!}\Sym^n_g(g')^{\ast}L\Gamma^n_{\ext} \ar@{-}[r]^\sim& R\Sym^n_S(u')_{!}L\Gamma^n_{\ext}g'^{\ast}.
}
$$

\item[${\rm(vi'')}$] 对每个\emph{根}基变换 $g:Y'\rightarrow Y$, 它给出 $S$-概形的笛卡儿图
$$
\xymatrix@C=2.5cm@R=1.5cm{
X' \ar[r]^{g'}\ar[d]^{u'}& X \ar[d]^{u}\\
Y' \ar[r]^{g}& Y,
}
$$
下列函子图交换:
\[
\xymatrix{
\Sym^n_S(g)^{\ast}L\Gamma^n_{\ext}Ru_{!} \ar@{-}[r]^{\sim} \ar[d]^{k^n_u}&
L\Gamma^n_{\ext}g^{\ast}Ru_{!} \ar[r]^{\sim}& L\Gamma^n_{\ext}
Ru'_{!}g'^{\ast} \ar[d]^{k^n_{u'}}\\
\Sym^n_S(g)^{\ast}R\Sym^n_S(u\sisi{_!)}{)_!L\Gamma^n_{\ext}} \ar[r]^{\sim}& R
\Sym^n_S(u')_{!}{g'}^{\ast}L\Gamma^n_{\ext} \ar@{-}[r]^{\sim}& R\Sym^n_S(u')_{!}L\Gamma^n_{\ext}g'^{\ast}
}
\]
\end{itemize}

\begin{enonce*}{定理 5.5.21}\etiquette[5.5.21]{XVII.thm5.5.21}
{\rm(对称\pageoriginale Künneth 公式).} 设 $S$ 为凝聚概形 (= 拟紧且拟分离), 以交换挠环层 $\Aa$ 赋环, 且其茎均为 Noether 环. 以 $\Aa$ 的逆像给每个 $S$-概形赋环. 设 $u:X\rightarrow Y$ 为拟射影 $S$-概形之间的态射. 对每个\sisi{}{ }整数 $n\geq0$ 以及每个 $K\in\Ob D^{b,\tor\leq0}(X,\Aa)$, 对称 Künneth 箭头 \eqhyperref[XVII.eq5.5.17.2]{5.5.17.2}
$$
k^n:L\Gamma^n_{\ext}Ru_{!}K\longrightarrow
R\Sym^n_S(u)_{!}L\Gamma^n_{\ext}K
$$
是 $D^{b,\tor\leq0}(\Sym^n_S(Y),\Aa)$ 中的同构.
\end{enonce*}

这一定理的证明将占据 \hyperref[XVII.subsec5.5]{5.5} 的其余部分.

\begin{enonce*}[remark]{注 5.5.21.1}\etiquette[5.5.21.1]{XVII.rem5.5.21.1}
拟射影性假设的作用仅在于保证 $u$ 容许紧化 $\bar u$
$$
X\xrightarrow{j}\bar X\xrightarrow{\bar u}Y
$$
使概形 $\Sym^n_S(X)$, $\Sym^n_S(\bar X)$ 及 $\Sym^n_S(Y)$ 对每个 $n$ 都存在. 代数空间的语言应能使我们将这一假设换成仅假设 $u$ 可紧化.
\end{enonce*}

证明的\pageoriginale 第一部分是通过分解归约（dévissage）, 将 \hyperref[XVII.thm5.5.21]{5.5.21} 化为一个“不可约”的情形:

\begin{enonce*}{引理 5.5.22}\etiquette[5.5.22]{XVII.lem5.5.22}
定理 \hyperref[XVII.thm5.5.21]{5.5.21} 来自下述特例:
\begin{itemize}
\item[${(\ast)}$] 对每个代数闭域 $k$, $k$ 上的每条光滑射影曲线 $X$, 以及每个素数 $\ell$, 对称 Künneth 箭头 (取 $Y=S=\Spec(k)$, $\Aa=\ZZ/\ell$, $K=\ZZ/\ell$ 所得; 见下文说明)
$$\etiquette[5.5.22.1]{XVII.eq5.5.22.1}
L\Gamma^n(R\Gamma(X,\ZZ/\ell))\longrightarrow
R\Gamma(\Sym^n_k(X),\ZZ/\ell)\leqno(5.5.22.1)
$$
都为同构 $(n\geq0)$.
\end{itemize}
\end{enonce*}

在 \sisi{\hyperref[XVII.eq5.5.22.1]{5.5.22.1}}{\eqhyperref[XVII.eq5.5.22.1]{5.5.22.1}} 中, 符号 $\Gamma$ 不幸有双重用途: 在 $R\Gamma$ 中它表示“整体截面”函子, 而在 $L\Gamma^n$ 中则表示 \hyperref[XVII.par5.5.2.3]{5.5.2.3} 的函子. 这一公式使用了同构 $\Gamma^n_{\ext}(\ZZ/\ell)\simeq\ZZ/\ell$.

设 $S$, $X$, $Y$, $u$ 及 $\Aa$ 如 \hyperref[XVII.thm5.5.21]{5.5.21} 中所述.

\begin{enonce*}{引理 5.5.22.2}\etiquette[5.5.22.2]{XVII.lem5.5.22.2}
设 $C^{b,\geq0}(X,\Aa)$ 中有 \ttor-维数 $\leq0$ 的复形的正合列:
$$
0\longrightarrow K'\longrightarrow K\longrightarrow K''\longrightarrow0{\quoi}.
$$
若关于复形 $K'$, $K$ 及 $K''$ 中任意两个的所有 Künneth 箭头 $k^n$ $(n\geq0)$ 均为同构, 则关于第三个的 Künneth 箭头也都为同构.
\end{enonce*}

我们\pageoriginale 化为 $u$ 射影、且复形 $K'$, $K$ 及 $K''$ 满足 \hyperref[XVII.lem5.5.16]{5.5.16} 的条件 (i) 与 (ii) 的情形, 其中 $N$ 至少为 $u$ 的纤维维数的两倍. 对称 Künneth 箭头此时由 \eqhyperref[XVII.eq5.5.17.1]{5.5.17.1} 给出.

对 $n$ 归纳证明: 若对于 $i\leq n$ 以及复形 $K'$, $K$ 及 $K''$ 中的任意两个, 箭头 $k^i$ 都为同构, 则第三个复形的箭头 $k^i$ $(i\leq n)$ 也都为同构.

为此注意, 由 \sisi{\hyperref[XVII.subsubsec5.5.14]{5.5.14}}{\hyperref[XVII.subsubsec5.5.15]{5.5.15}}, 态射
\begin{equation*}
k^n:L\Gamma^n_{\ext}Ru_{\ast}K\longrightarrow
R\Sym^n(u)_{\ast}L\Gamma^n_{\ext}K\tag{a}
\end{equation*}
可以表示为有限滤过的滤过复形之间的态射
$$
{k'}^n:K_0\longrightarrow K_1
$$
在逐次商上, 这一态射诱导的箭头同构于复合箭头
\begin{equation*}
L\Gamma^i_{\ext}Ru_{\ast}K'\overset{\LL}{\vee}
L\Gamma^{n-i}_{\ext}Ru_{\ast}K''\longrightarrow
R\Sym^n(u)_{\ast}(L\Gamma^i_{\ext}K'\overset{\sisi{}{\LL}}{\vee}
L\Gamma^{n-i}_{\ext}K'')\tag{$\text{b}_i$}
\end{equation*}
具体如下:
\begin{gather*}
L\Gamma^i_{\ext}Ru_{\ast}K'\overset{\LL}{\vee}
L\Gamma^{n-i}_{\ext}Ru_{\ast}K''\xrightarrow{k^i\overset{\LL}{\vee}k^{n-i}}
R\Sym^i(u)_{\ast}L\Gamma^i_{\ext}K'\overset{\LL}{\vee}
R\Sym^{n-i}(u)_{\ast}L\Gamma^{n-i}_{\ext}(K'')\\
\underset{\defn}{=}
\vee_{i,n-i{\ast}}(R\Sym^i(u)_{\ast}L\Gamma^i_{\ext}K'\overset{\LL}{\boxtimes}
R\Sym^{n-i}(u)_{\ast}L\Gamma^{n-i}_{\ext}(K''))\\
\xrightarrow{\substack{\hyperref[XVII.thm5.4.3]{5.4.3}\\\sim}}
\vee_{i,n-i{\ast}}\sisi{(R\Sym^i(u)_{\ast}\times R\Sym^{n-i}(u)_{\ast})}{(R\Sym^i(u)\times R\Sym^{n-i}(u))_\ast}
(L\Gamma^i_{\ext}K'\overset{\LL}{\boxtimes}
L\Gamma^{n-i}_{\ext}K'')\\
=R\Sym^n(u)_{\ast}\vee_{i,n-i{\ast}}(L\Gamma^i_{\ext}K'\overset{\LL}{\boxtimes}L\Gamma^{n-i}_{\ext}K'')
\overset{\dfn}{=}R\Sym^n(u)_{\ast}(L\Gamma^i_{\ext}K'\overset{\LL}{\vee}L\Gamma^{n-i}_{\ext}K''){\quoi}.
\end{gather*}

由\pageoriginale 归纳假设 (及 Künneth 公式), 这些箭头在 $i\neq0,n$ 时为拟同构. 因而下列三个断言中的任意两个蕴含第三个: (a) 为拟同构; $(\text{b}_0)$ 为拟同构; $(\text{b}_n)$ 为拟同构. 这些恰是所考察的三条箭头, \cqfd

\setcounter{subsubsection}{22}
\setcounter{paragraph}{2}
\paragraph{}\etiquette[5.5.22.3]{XVII.par5.5.22.3}
假设 \hyperref[XVII.lem5.5.22]{5.5.22} 的断言 $(\ast)$, 来证明 \hyperref[XVII.thm5.5.21]{5.5.21}. 因而设 $S$, $X$, $Y$, $\Aa$, $K$ 及 $u$ 如 \hyperref[XVII.thm5.5.21]{5.5.21}. 问题在 $S$ 上是局部的, 可设它仿射. 每个复形 $K$ 都是形如 $i_{!}i^{\ast}K$ 的复形的逐次扩张, 其中 $i$ 为局部闭仿射子集到 $X$ 的包含映射. 通过分解归约 (\hyperref[XVII.lem5.5.22.2]{5.5.22.2}), 可假设 $K$ 形如 $i_{!}K_1$, 其中 $i:X_1\rightarrow X$ 为起点仿射的浸入. 相容性 \hyperref[XVII.par5.5.20.2]{5.5.20.2} 允许用 $X_1$ 代替 $X$, 用 $K_1$ 代替 $K$, 因而可设 $u$ 仿射. 将 $u$ 分解为纤维维数 $\leq1$ 的态射的复合 $u=u_1\ldots u_k$. 相容性 \hyperref[XVII.par5.5.20.2]{5.5.20.2} 表明, 只须证明关于 $u_i$ 及任意复形 $K$ 的箭头 $k^n$ 为同构. 因而问题化为 $S$ \emph{仿射且 $u$ 的相对维数为 $\leq1$}的情形.

\paragraph{}\etiquette[5.5.22.4]{XVII.par5.5.22.4}
要使 $k^n$ 为同构, 只须它在 $\Sym^n_S(Y)$ 的几何纤维上的限制为同构; 因而相容性 \hyperref[XVII.par5.5.20.3]{5.5.20.3} ${\rm(v'')}$ 将问题化为 $S$ 为\emph{代数闭域 $k$ 的谱}的情形.

\paragraph{}\etiquette[5.5.22.5]{XVII.par5.5.22.5}
对\pageoriginale $\Sym^n_k(Y)$ 的每个有理点 $y$, 存在 $Y$ 的一个有限既约子概形 $Y'$, 使 $\sisi{x}{y}$ 属于 $\Sym^n_k(Y')$ 的像. 因而相容性 \hyperref[XVII.par5.5.20.3]{5.5.20.3} ${\rm(vi'')}$ 允许假设 $Y$ \emph{为有限个} $\Spec(k)$ \emph{的副本的不交并}.

\paragraph{}\etiquette[5.5.22.6]{XVII.par5.5.22.6}
若 $Y=\amalg Y_i$, 则 $X$ 上的每个复形 $K$ 都是支集在某个概形 $X_i=u^{-1}(Y_i)$ 上的复形之和. 由分解归约 (\hyperref[XVII.lem5.5.22.2]{5.5.22.2}) 以及相容性 \hyperref[XVII.par5.5.20.3]{5.5.20.3} ${\rm(vi'')}$, 只须对每个态射 $u_i:X_i\rightarrow Y_i$ 证明 \hyperref[XVII.thm5.5.21]{5.5.21}. 根据 \hyperref[XVII.par5.5.22.5]{5.5.22.5}, 这将问题化为
$$
\text{\rm $Y=S=\Spec(k)$, $k$ 代数闭, $\dim X\leq1$.}
$$

层 $\Aa$ 此时化为一个 Noether 环 $A$. 下述命题使我们能化为 $A$ 为域的情形.

\begin{enonce*}{命题 5.5.23}\etiquette[5.5.23]{XVII.prop5.5.23}
设 $A$ 为 (交换) Noether 环, $K\in\Ob D^{-}(A)$ 为 $A$-模的上有界复形, $n$ 为整数. 下列条件等价:

\begin{enumerate}
\renewcommand{\theenumi}{\roman{enumi}}
\renewcommand{\labelenumi}{(\theenumi)}
\item $K$ 的 \tTor-维数为 $\leq n$.

\item 对每个 $x\in \Spec (A)$ 以及每个 $i<-n$, 都有 $H^{i}(K\overset{\LL}{\underset{A}{\otimes}} k(x))=0$.
\end{enumerate}
\end{enonce*}

显然有 (i) $\Rightarrow$ (ii)\sisi{,}{.} 假设 (ii). 为证明 (i), 对 $\Spec(A)$ 的闭子集 $F$ 作 Noether 归纳, 证明若模 $N$ 在 $x\not\in F$ 时满足 $N_{x}=0$, 则在 $i<-n$ 时有 $H^{i}(K\overset{\LL}{\otimes}_{A}N)=0$.

设\pageoriginale $I$ 为定义 $F$ 且使 $A/I$ 既约的唯一理想, $N_{k}$ 为 $N$ 中被 $I^{k}$ 零化的子模. 有 $N=\varinjlim N_{k}$. 由于 Tor 函子与滤过归纳极限可交换, 可化为存在 $k$ 使 $N=N_{k}$ 的情形, 从而 $N$ 是被 $I$ 零化的模的逐次扩张; 上同调长正合列将问题化为 $N$ 本身已是 $A/I$-模的情形.

设 $\eta$ 为 $F$ 的一个极大点, 并由正合列
$$\etiquette[5.5.23.1]{XVII.eq5.5.23.1}
0\longrightarrow M_{1}\longrightarrow N \xrightarrow{i} N\otimes
k(\eta)\longrightarrow M_{2} \longrightarrow 0{\quoi}.\leqno(5.5.23.1)
$$
定义 $A/I$-模 $M_{1}$ 与 $M_{2}$.

有 $(M_{j})_{\eta}=0$ $(j=1,2)$. 对每个满足 $\eta\not\in G$ 的闭子集 $G\subset F$, 若以 $M_{j,G}$ 表示 $M_{j}$ 中由 $M_{j}$ 的支集在 $G$ 中的截面组成的子模, 则
$$
M_{j}=\varinjlim_{\eta\not\in G} M_{j,G}{\quoi}.
$$
由归纳假设及极限过渡, 因而有
$$
H^{i}(K\overset{\LL}{\otimes}_{A} M_j)=0\quad\mbox{\rm
当}\quad i< -n{\quoi}.
$$

模 $N\otimes k(\eta)$ 为 $k(\eta)$ 上的向量空间, 因而为若干个 $k(\eta)$ 的直和. 由假设, 有
$$
H^{i}(K\overset{\LL}{\otimes} (N\otimes
k(\eta))\sisi{}{)}=0\quad\mbox{\rm 当}\quad i<-n{\quoi}.
$$

若 $I$ 为 \eqhyperref[XVII.eq5.5.23.1]{5.5.23.1} 中箭头 $i$ 的像, 则有正合列
\begin{align*}
& H^{i-1}(K\overset{\LL}{\otimes}M_{2})\longrightarrow
H^{i}(K\overset{\LL}{\otimes}I)\longrightarrow
H^{i}(K\overset{\LL}{\otimes}(N\otimes k(\eta)))\\
& H^{i}(K\overset{\LL}{\otimes}M_{1})\longrightarrow
H^{i}(K\overset{\LL}{\otimes}N)\longrightarrow
H^{i}(K\overset{\LL}{\otimes}I){\quoi},
\end{align*}
从而\pageoriginale 在 $i<-n$ 时有 $H^{i}(K\overset{\LL}{\otimes}N)=0$,
\cqfd

\begin{enonce*}{推论 5.5.24}\etiquette[5.5.24]{XVII.cor5.5.24}
要使 $D^{-}(A)$ 中的态射 $u:K\rightarrow L$ 为同构, 只须对每个 $x\in \Spec(A)$, 态射
$$
u_{x}:K\overset{\LL}{\otimes}_{A}k(x)\longrightarrow
L\overset{\LL}{\otimes}_{A}k(x)
$$
都是 $D^{-}(k(x))$ 中的同构.
\end{enonce*}

\sisi{\noindent}{}设 $M$ 为 $u$ 的锥 (映射柱). 由假设, 对每个 $x\in\Spec(A)$, 都有 $M\overset{\LL}{\otimes}_{A}k(x)=0$. 因而由 \hyperref[XVII.prop5.5.23]{5.5.23}, 有 $M=0$, 所以 $u$ 为同构.

\setcounter{subsubsection}{25}
\setcounter{paragraph}{0}
\paragraph{}\etiquette[5.5.25.1]{XVII.par5.5.25.1}
继续在 \hyperref[XVII.prop5.5.23]{5.5.23} 处中断的 \hyperref[XVII.lem5.5.22]{5.5.22} 的证明. 由相容性 \hyperref[XVII.par5.5.20.3]{5.5.20.3} ${\rm(iv'')}$ 以及 \hyperref[XVII.cor5.5.24]{5.5.24}, 与 $K\in \Ob D^{b,\tor\leq 0}(X,A)$ 相伴的态射 $k^{n}$ 为同构, 当且仅当对 $x\in
\Spec(A)$, 与 $K\overset{\LL}{\otimes}_{A}k(x)\in \Ob
D^{b,\tor\leq 0}(X,k(x))$ 相伴的态射 $k^{n}$ 为同构. 这将问题化为 $A$ \emph{为特征} $\ell > 0$ \emph{的域}的情形.

\paragraph{}\etiquette[5.5.25.2]{XVII.par5.5.25.2}
正合列 (其中 $C$ 表示一个锥)
$$
0\longrightarrow K[-1]\longrightarrow CK[-1]\longrightarrow
K\longrightarrow 0
$$
表明, 对 $K$ 的态射 $k^{n}$ 为同构, 当且仅当对 $K[-1]$ 的相应态射为同构 (\hyperref[XVII.lem5.5.22.2]{5.5.22.2}).

\paragraph{}\etiquette[5.5.25.3]{XVII.par5.5.25.3}
正合列\pageoriginale
$$
0\longrightarrow \sigma_{\geq n} K\longrightarrow K\longrightarrow
\sigma_{\leq n-1} K\longrightarrow 0
$$
表明, 要使 $K$ 为平坦复形时的 Künneth 态射为同构, 只须它们对 $K$ 的各分量放在次数 $0$ 处时为同构 (\hyperref[XVII.lem5.5.22.2]{5.5.22.2}\sisi{)}{} 及 \sisi{(}{}\hyperref[XVII.par5.5.25.2]{5.5.25.2}).

\paragraph{}\etiquette[5.5.25.4]{XVII.par5.5.25.4}
于是设 $\sisi{F}{K}$ 为 $A$-向量空间层. 这一层为可构造层的滤过归纳极限 (\hyperrefext[IX][IX.cor2.7.2]{2.7.2}); 由于 $A$ 为域, 这些层自动平坦; 通过极限过渡, 问题化为 $K$ \emph{可构造}且集中于零次的情形.

\paragraph{}\etiquette[5.5.25.5]{XVII.par5.5.25.5}
若复形 $K$ 为复形 $L$ 的直和因子, 则与 $K$ 相伴的态射 $k^{n}$ 为与 $L$ 相伴的态射 $k^{n}$ 的直和因子, 因而后者为同构时前者也为同构. 所以迹方法 (\hyperrefext[IX][IX.rem5.8]{5.8}) 适用, 只须处理 $F$ 形如 $j_{!}A_{X'}$ 且 $j:X'\rightarrow X$ 拟有限的情形. 相容性 \hyperref[XVII.par5.5.20.2]{5.5.20.2} 允许用 $X'$ 代替 $X$, 从而得到归约
$$
\mbox{$X$ \emph{的维数为 $\leq 1$, 且} $K=A_{X}${\quoi}.}
$$

\paragraph{}\etiquette[5.5.25.6]{XVII.par5.5.25.6}
可假设 $X$ 既约. 若 $j:U\rightarrow X$ 为 $X$ 的光滑部分, 补集为 $i:F\rightarrow X$, 则正合列
$$
0\longrightarrow j_{!}j^{\ast}A_{X}\longrightarrow
A_{X}\longrightarrow i_{!}i^{\ast}A_{X}\longrightarrow 0
$$
通过分解归约, 将问题化为 $X$ \emph{或为光滑曲线, 或为有限集}的情形. 此外, 由公式
$$
A_{X}=A\overset{\LL}{\otimes}_{\ZZ/\ell}(\ZZ/\ell)_{X}
$$
及 \hyperref[XVII.par5.5.20.3]{5.5.20.3} ${\rm(iv'')}$, 问题化为 $A_{X}=\ZZ/\ell$ 的情形.

若\pageoriginale $X$ 为曲线, 则将 $X$ 嵌入相应的光滑射影曲线, 再作一次分解归约, 即化为下列两种情形之一:
\begin{enumerate}
\renewcommand{\theenumi}{\roman{enumi}}
\renewcommand{\labelenumi}{(\theenumi)}
\item $X$ 为光滑射影曲线, $Y=S=\Spec(k)$ 且 $k$ 代数闭, $\Aa=\ZZ/\ell$, $K=\ZZ/\ell$.

\item $X$ 为有限个 $\Spec(k)$ 的副本的不交并, $Y=S=\Spec(k)$ ($k$ 代数闭), $\Aa=\ZZ/\ell$, $K=\ZZ/\ell$.
\end{enumerate}

最后将 $\sisi{}{(}\ZZ/\ell\sisi{}{)}_{X}$ 分解为层 $\sisi{}{(}\ZZ/\ell\sisi{}{)}_{X\sisi{}{'}}$, 其中 $X'$ 为 $X$ 的连通分支, 就可在 (i) 及 (ii) 中假设 $X$ 连通. 此时情形 (ii) 是平凡的. 这就证明了 \hyperref[XVII.lem5.5.22]{5.5.22}.

\begin{enonce*}{引理 5.5.26}\etiquette[5.5.26]{XVII.lem5.5.26}
定理\emph{\hyperref[XVII.thm5.5.21]{5.5.21}} 来自\emph{\hyperref[XVII.lem5.5.22]{5.5.22}} 的断言 $(\ast)$ 的下列两个特例:
\begin{itemize}
\item[\emph{情形} 1.] 假设 $k=\CC$.

\item[\emph{情形} 2.] 假设 $k$ 的特征为 $p>0$ 且 $\ell=p$.
\end{itemize}
\end{enonce*}

用 \hyperref[XVII.lem5.5.22]{5.5.22} 的记号, 以 $p$ 表示 $k$ 的特征. 若 $p=0$, 则 Lefschetz 原理 (这里依赖于相容性 \hyperref[XVII.par5.5.20.3]{5.5.20.3} ${\rm(v'')}$ 以及 XII \hyperref[XII.cor5.4]{5.4}) 允许化为 $k=\CC$ 的情形.

只剩处理 $p>0$ 且 $\ell\neq p$ 的情形. 曲线 $X$ 可提升为 $k$ 上的 Witt 向量环 $W(k)$ 上的一条曲线 $X'$ (\SGA 1 III 7.4). 令 $S=\Spec(W(k))$, 并考虑图
$$
\xymatrix@C=2.25cm@R=1.5cm{
X'_{\eta} \ar[d]^{f_\eta}          \ar[r]              & X' \ar[d]^{f'}& X\ar[l]\ar[d]^{f}\\
\eta               \ar[r]         & S   & \Spec(k), \ar[l],
}
$$
并\pageoriginale 设 $f_{n}$, $f'_{n}$ 及 $f_{\eta n}$ 分别为 $\Sym^{n}_{k}(X)$, $\Sym^{n}_{S}(X')$ 及 $\Sym^{n}_{\eta}(X'_{\eta})$ 到 $\Spec(k)$, $S$ 及 $\eta$ 的投影. 已知 $f'_{n}$ 为光滑射影态射 (射影性显然; 对光滑性, 最简单的方法是化为 $X'=\PP^{1}_{S}$ 的情形, 此时由初等对称函数定理有 $\Sym^{n}_{S}(X')\simeq
\PP^{n}$). 由特殊化定理 \hyperrefext[XVI][XVI.cor2.2]{2.2}, 层 $R^{i}f'_{n{\ast}}$ $(\ZZ/\ell)$ 在 $S$ 上局部常值. 层 $L^{i}\Gamma^{n}(Rf'_{\ast}\ZZ/\ell)$ 也如此. 因而, 要验证与 $k$ 上的 $X$ 相伴的 \eqhyperref[XVII.eq5.5.22.1]{5.5.22.1} 中的态射为同构, 由 \hyperref[XVII.par5.5.20.3]{5.5.20.3} ${\rm(v'')}$, 只须证明 $X$ 在 $S$ 上的几何泛纤维 $X_{\bar{\eta}}$ 的类似陈述. 域 $k(\bar{\eta})$ 的特征为 $0$, 因而根据第一次归约, 这就证明了 \hyperref[XVII.lem5.5.26]{5.5.26}.

为在 \hyperref[XVII.lem5.5.26]{5.5.26} 的两个特例中证明 \hyperref[XVII.thm5.5.21]{5.5.21}, 我们将对称 Künneth 态射一方面与其超越类似物比较 (\hyperref[XVII.subsubsec5.5.27]{5.5.27} 及 \hyperref[XVII.subsubsec5.5.28]{5.5.28}), 另一方面与其凝聚类似物比较 (\hyperref[XVII.subsubsec5.5.29]{5.5.29} 至 \hyperref[XVII.lem5.5.37]{5.5.37}).

\setcounter{subsubsection}{26}
\subsubsection{}\etiquette[5.5.27]{XVII.subsubsec5.5.27}
若 $p:X\rightarrow S$ 为拓扑空间之间的连续映射, $S$ 以交换环层 $\Aa$ 赋环, 则 \hyperref[XVII.subsubsec5.5.7]{5.5.7} 及 \hyperref[XVII.subsubsec5.5.8]{5.5.8} 的构造原样适用, 定义从 $X$ 上的模到 $\Sym^{n}_{S}(X)=X^{n}/\Sf_{n}$ 上的模的函子 $\Gamma^{n}_{\ext}$. \hyperref[XVII.subsubsec5.5.7]{5.5.7} 至 \hyperref[XVII.prop5.5.18]{5.5.18} 的形式论证均可照搬. 同样, 若 $S$, $X$ 及 $Y$ 为局部紧 Hausdorff 拓扑空间,
$$
\xymatrix@C=1.5cm@R=1.25cm{
X \ar[rr]^{u}\ar[rd]&  &  Y\ar[ld]\\
            & S &,
}
$$
并且\pageoriginale $u$ 的纤维的上同调维数有界, 则 \hyperref[XVII.subsubsec5.5.19]{5.5.19} 及 \hyperref[XVII.subsubsec5.5.20]{5.5.20} 中关于固有支集上同调中的对称 Künneth 态射的形式体系也均可照搬.

若 $p:X\rightarrow S$ 为 $\CC$ 上有限型概形之间的态射, 则函子 $\Gamma^{n}$ 与 $L\Gamma^{n}$ 同“过渡到超越情形”的函子可交换. 最后, 若 $X$, $Y$, $S$, $\Aa$ 及 $K$ 如 \hyperref[XVII.thm5.5.21]{5.5.21}, $S$ 在 $\CC$ 上分离且有限型, 并且以 $\epsilon$ 统称 $\epsilon:X_{\cl}\rightarrow
X_{\et}$ 类型的景态射, 则下图交换:
$$
\vcenter{\etiquette[5.5.27.1]{XVII.eq5.5.27.1}
\xymatrix@C=1.75cm@R=1.5cm{
\epsilon^{\ast} L\Gamma^{n}_{\ext} Ru_{!}K  \ar[r]^{k^{\sisi{k}{n}}} \ar[d]&
\epsilon^{\ast} R \Sym^{n}_{S}(u)_{!} L \Gamma^{n}_{\ext} K \ar[d]\\
L\Gamma^{n}_{\ext} R u^{\an}_{!} \epsilon^{\ast}K  \ar[r]^{k^{n \an}}
& R \Sym^{n}_{S}(u^{\an})_{!} L\Gamma^{n}_{\ext} \epsilon^{\ast}K.
}}\leqno{(5.5.27.1)}
$$

\subsubsection{}\etiquette[5.5.28]{XVII.subsubsec5.5.28}
由相容性 \eqhyperref[XVII.eq5.5.27.1]{5.5.27.1} 以及比较定理的一个初等情形 (\hyperrefext[XI][XI.thm4.4]{4.4} 或 \hyperrefext[XVI][XVI.thm4.1]{4.1}), 要处理 \hyperref[XVII.lem5.5.26]{5.5.26} 的情形 1, 只须验证: 若 $X$ 为光滑复射影曲线, 则 \eqhyperref[XVII.eq5.5.22.1]{5.5.22.1} 的拓扑类似物为同构.

拓扑空间 $X^{\an}$ 可三角剖分 (这一事实在此是平凡的, 因为 $X$ 为 Riemann 球面的分歧覆盖). 分解归约引理 \hyperref[XVII.lem5.5.22.2]{5.5.22.2} 的拓扑类似物, 于是将问题化为证明
\begin{itemize}
\item[$(\ast)$] 若\pageoriginale $D$ 为闭单形 (维数为 0、1 或 2), 则拓扑的对称 Künneth 态射
$$\etiquette[5.5.28.1]{XVII.eq5.5.28.1}
k^{n}:L\Gamma^{n} R\Gamma(D,\ZZ/\ell)\longrightarrow R\Gamma(\Sym^{n}(D\sisi{}{)},\ZZ/\ell)\leqno(5.5.28.1)
$$
为同构.
\end{itemize}

有
$$
R\Gamma(D,\ZZ/\ell)=\ZZ/\ell\;\mbox{\rm 置于次数 }
0{\quoi},
$$
因而
$$
L\Gamma^{n}R\Gamma(D,\ZZ/\ell)=\ZZ/\ell\; \mbox{\rm
置于次数 $0$}{\quoi}.
$$

另一方面, $\Sym^{n}(D)$ 可缩, 于是仍有
$$
R\Gamma(\Sym^{n}(D), \ZZ/\ell)=\ZZ/\ell\;\mbox{\rm 置于
次数 $0$}{\quoi}.
$$

因而 \eqhyperref[XVII.eq5.5.28.1]{5.5.28.1} 两边的 $H^{i}$ 在非零次数都为零, 而在 $H^0$ 上诱导的态射等同于 $\ZZ/\ell$ 的恒等态射.

\subsubsection{}\etiquette[5.5.29]{XVII.subsubsec5.5.29}
对于我们打算作的应用, 将概形 $X$ 上的拟凝聚层视为 $X$ 的平展景上的拟凝聚 $\Oo$-模层是有用的.

设 $p:X\rightarrow \sisi{X}{S}$ 为概形的拟射影态射 (参见 \hyperref[XVII.rem5.5.21.1]{5.5.21.1} \footnote[1]{以下直至包括 \hyperref[XVII.prop5.5.34]{5.5.34} 的全部内容, 在仅假设 $X$ 与 $Y$ 容许 $S$ 上的相对丰沛可逆层的条件下仍然成立 (不假设 $X$ 与 $Y$ 在 $S$ 上有限型).}). 若 $ℱ$ 为 $X_{\text{\rm et}}$ 上的 $\Oo_{X}$-模层, 则
\begin{enumerate}
\renewcommand{\theenumi}{\alph{enumi}}
\renewcommand{\labelenumi}{\theenumi)}
\item 以 $\bigotimes\limits^{n}ℱ$ 表示 $(X/S)^{n}_{\text{\rm et}}$ 上的模层, 它由 $ℱ$ 在 $(X/S)^{n}_{\et}$ 上的逆像 $\Oo_{(X/S)^n}$-模的张量积给出;

\item 层 $\bigotimes\limits^{n}ℱ$ 为 $\mathcal{G}_n$-等变层. 取不变量后, 它定义 $\Sym^{n}_{S}(X)$ 上的一个层 $TS^{n}_{\text{\rm ext}}(ℱ)$: 若 $\sigma$ 为 $(X/S)^{n}$ 到 $\Sym^{n}_{S}(X)$ 的投影, 则有
$$\etiquette[5.5.29.1]{XVII.eq5.5.29.1}
TS^{n}_{\ext}(ℱ)=(\sigma_{\ast}\bigotimes^{n}ℱ)^{\Sf^{n}}{\quoi}.\leqno(5.5.29.1)
$$
\end{enumerate}

\emph{假设\pageoriginale $p$ 平坦}; 如 \hyperref[XVII.subsubsec5.5.8]{5.5.8} 中那样, 将函子 $\Gamma^{n}_{\text{\rm ext}}$ 定义为 $TS^{n}_{\ext}$ 的第 $0$\sisi{阶}{阶}左导出函子. 用 \hyperref[XVII.par5.5.2.5]{5.5.2.5} 及 \hyperref[XVII.subsubsec5.5.8]{5.5.8} 的记号, 若
$$
P_{1}\longrightarrow P_{0}\longrightarrow ℱ\longrightarrow 0
$$
为 $ℱ$ 的平坦表示, 则序列
$$
\xymatrix@C=1.75cm@M=.25cm{
TS^{n}_{\ext} (P'_1) \ar@<3pt>[r]^{j_0} \ar@<-3pt>[r]_{j_1}& TS^{n}_{\ext}(P_0)\ar[r] & \Gamma^{n}_{\ext}(\sisi{F}{ℱ})
}
$$
正合. 若 $x$ 为如 \hyperref[XVII.par5.5.7.1]{5.5.7.1} 中所述的 $\Sym^{n}_{S}(X)$ 的几何点, 则有 (\cf \eqhyperref[XVII.eq5.5.8.1]{5.5.8.1})
\begin{align*}
& \Gamma^{n}_{\text{\rm
ext}}(ℱ)_{x}=\Oo_{x}\otimes_{C}(\bigotimes^{\ell}_{1}
\Gamma_{\Oo_{S}}^{n_{i}}(\sisi{F}{ℱ}_{y_{i}})){\quoi},\tag{5.5.29.2}\\
& C=\bigotimes^{\ell}_{1} \Sym^{n_{i}}_{\Oo_{S}}(\Oo_{y_{i}}){\quoi}.\tag*{其中}
\end{align*}\etiquette[5.5.29.2]{XVII.eq5.5.29.2}% 删除此注释会引入不需要的空白.
事实上, 只须在 $\sisi{F}{ℱ}$ 的茎为自由模的平凡情形下验证这一公式. 注意, $\Oo_{x}$ 在 $C$ 上平坦, 因为它是 $C$ 在 $x$ 的像这一闭点处的严格 Hensel 化.

\subsubsection{}\etiquette[5.5.30]{XVII.subsubsec5.5.30}
假设 $X$ 与 $S$ 仿射, $p$ 平坦且 $ℱ$ 拟凝聚. 设 $S=\Spec(A)$, $X=\Spec(B)$, 并令 $F=H^0(X,ℱ)$, 则 $ℱ$ 为 $B$-模 $F$ 定义的层 $F^{\sim}$. 由 \eqhyperref[XVII.eq5.5.29.1]{5.5.29.1}, 有
$$
TS^{n}_{\ext}(ℱ)=TS^{n}_{A}(F)^{\sim}{\quoi},
$$
因而导出后有
$$\etiquette[5.5.30.1]{XVII.eq5.5.30.1}
\Gamma^{n}_{\text{\rm
ext}}(ℱ)=\Gamma^{n}_{A}(F)^{\sim}{\quoi}.\leqno(5.5.30.1)
$$
所以\pageoriginale 函子 $\Gamma^{n}_{\text{\rm ext}}$ 表现为一个“相对于 $A$”的函子 $\Gamma^{n}$, 正如 \eqhyperref[XVII.eq5.5.29.2]{5.5.29.2} 已在局部所表述的.

\subsubsection{}\etiquette[5.5.31]{XVII.subsubsec5.5.31}
以 $D^{b,\tor_{S}\leq 0}(X,\Oo_{X})$ 表示 $D^{b}(X,\Oo_{X})$ 中由 $\Oo_{X}$-模复形组成的子范畴, 这些复形作为 $p^{-1}(\Oo_{S})$-模复形的 \tTor-维数为 $\leq 0$.

考虑条件

\paragraph{}\etiquette[5.5.31.1]{XVII.par5.5.31.1}
$K\in \Ob C^{b,\geq 0}(X,\sisi{O}{\Oo}_{X})$, 并且在 $X$ 的每个几何点 $x$ 处, $K_{x}$ 同伦于一个非负次数的 $\Oo_{X,x}$-模复形, 其分量在 $\Oo_{S,p(x)}$ 上平坦.

由于态射 $p$ 平坦, 从而有“足够多”的在 $S$ 上平坦的 $\Oo_{X}$-模, 如 \hyperref[XVII.lem5.5.13.2]{5.5.13.2} 中一样可验证, $D^{b,\tor_{S}\leq 0}(X,\Oo_{X})$ 是从满足 (\hyperref[XVII.par5.5.31.1]{5.5.31.1}) 的复形 $K\in \Ob K^{b,\geq 0}(X,\Oo_{X})$ 的范畴, 通过将拟同构变为可逆态射而得到的.

由 \eqhyperref[XVII.eq5.5.29.2]{5.5.29.2} 及 \hyperref[XVII.lem5.5.6.1]{5.5.6.1}, 若 $u:K\rightarrow L$ 为满足 (\hyperref[XVII.par5.5.31.1]{5.5.31.1}) 的复形之间的拟同构, 则 $\Gamma^{n}_{\text{\rm ext}}(u)$ (在 \textsc{Dold-Puppe} 的意义下) 仍为拟同构. 因而可以导出函子 $\Gamma^{n}_{\text{\rm ext}}$, 得到
$$\etiquette[5.5.31.2]{XVII.eq5.5.31.2}
L\Gamma^{n}_{\text{\rm ext}}:D^{b,\tor_{S}\leq
0}(X,\Oo_{X})\longrightarrow D^{b,\tor_{S}\leq 0}(\Sym^{n}_{S}(X), \Oo_{\Sym^{n}_{S}(X)})\leqno(5.5.31.2)
$$

\subsubsection{}\etiquette[5.5.32]{XVII.subsubsec5.5.32}
以 $D^{b,\tor_{S}\leq 0}_{\text{\rm qcoh}}(X,\Oo_{X})$ 表示 $D^{b,\tor_{S}\leq 0}(X,\Oo_{X})$ 中由上同调拟凝聚的复形组成的子范畴.

标准的\pageoriginale 论证使我们可以如 \hyperref[XVII.subsubsec5.5.17]{5.5.17} 中那样定义如下类型的对称 Künneth 箭头:

设 $u:X\rightarrow Y$ 为在 $S$ 上平坦的拟射影 $S$-概形之间的态射; 对 $K\in \Ob D^{b,\tor_{S}\leq 0}_{\text{\rm
qcoh}}(X,\Oo_{X})$, $D^{b,\tor_{S}\leq 0}_{\text{\rm qcoh}}(\Sym^{n}_{S}(Y),\Oo)$ 中的对称 Künneth 箭头 $k^{n}$ 为箭头
$$\etiquette[5.5.32.1]{XVII.eq5.5.32.1}
k^{n}:L\Gamma^{n}_{\text{\rm ext}}Ru_{\ast}K\longrightarrow
R\Sym^{n}_{S}u_{\ast}L\Gamma^{n}_{\text{\rm ext}} K{\quoi}.\leqno(5.5.32.1)
$$

\eqhyperref[XVII.eq5.5.32.1]{5.5.32.1} 将用到的形式性质的验证留给读者, 特别是如下性质:

\begin{enonce*}{引理 5.5.33}\etiquette[5.5.33]{XVII.lem5.5.33}
若 $S$ 的特征为 $p>0$, 且 $u:X\rightarrow S$ 固有且平坦, 则对称 Künneth 箭头 \eqhyperref[XVII.eq5.5.17.2]{5.5.17.2} 与 \eqhyperref[XVII.eq5.5.32.1]{5.5.32.1} 的图交换 ($Y=S, K=\sisi{\sheafZZ}{\ZZ}/p$ 或 $\Oo_{X}$):
\[
\xymatrix{
L\Gamma^{n}_{\sisi{\sheafZZ}{\ZZ}/p}Ru_{\ast}(\sisi{\sheafZZ}{\ZZ}/p)
\ar[r]^{k^n} \ar[d]&R
\Sym^{n}_S(u)_{\ast} (\sisi{\sheafZZ}{\ZZ}/p) \ar[d]&\\
L\Gamma^{n}_{\Oo} Ru_{\ast}(\Oo) \ar[r]^{k^n}& R
\Sym^{n}_{S}(u)_{\ast}\Oo &.
}
\]
\end{enonce*}

\begin{enonce*}{命题 5.5.34}\etiquette[5.5.34]{XVII.prop5.5.34}
在 \hyperref[XVII.subsubsec5.5.32]{5.5.32} 的假设下, 对称 Künneth 箭头 \eqhyperref[XVII.eq5.5.32.1]{5.5.32.1} 为同构.
\end{enonce*}

\begin{enumerate}
\renewcommand{\theenumi}{\alph{enumi}}
\renewcommand{\labelenumi}{\theenumi)}
\item 问题在 $S$ 上是局部的, 可以假设它仿射.

\item 对 $\Sym^{n}_{S}(Y)$ 的每个点 $x$, 存在 $Y$ 的一个仿射开集 $Y_{1}$, 使 $x$ 属于 $\Sym^{n}_{S}(Y_{1})$. 因而通过局部化可以假设 $Y$ 仿射.

\item 所考察类型的每个\pageoriginale 复形 $K$ 都可表示为 $X$ 上的一个拟凝聚模复形 $K_{0}$, 在次数 $<0$ 处为零. 若 $\mathcal{U}$ 为 $X$ 的一个仿射覆盖, 则 $K$ 容许由形如 $j_{\ast}j^{\ast}K$ 的复形的和组成的有限消解, 其中 $j$ 为仿射开集的包含映射 (交错 Čech 复形). 利用分解归约引理 \hyperref[XVII.lem5.5.22.2]{5.5.22.2} 的拟凝聚类似物, 可化为存在 $X$ 的仿射开集 $j:X_{1}\hookrightarrow X$ 及 $K_{1}\in D^{b,\tor_{S}\leq 0}_{\text{\rm
qcoh}}(X_{1},\Oo)$ 使 $K=Rj_{\ast}K_{1}$ 的情形.

\item 根据传递性 \hyperref[XVII.par5.5.20.2]{5.5.20.2} 的拟凝聚类似物, 只须对 $uj$ 与 $j$ 证明 \hyperref[XVII.prop5.5.34]{5.5.34}. 将归约 b) 应用于 $j$, 最终可化为 $X$, $Y$ 及 $S$ 仿射的情形. 此时可将 $K$ 表示为分量为 \sisi{}{$S$-}平坦拟凝聚模的复形 $K_{2}\in C^{b,\geq
0}(X,\Oo_{X})$ (\sisi{\quad}{\hyperref[XVII.def4.1.9]{4.1.9}}). 于是计算函子 $Ru_{\ast}$ 与 $L\Gamma^{n}$ 无须再消解 $K_{2}$, 只须注意, 由 \eqhyperref[XVII.eq5.5.30.1]{5.5.30.1}, 对 $X$ 上的拟凝聚模 $ℱ$, 有
\[
\Gamma^{n}_{\text{\rm ext}} u_{\ast}ℱ\isomto
\Sym^{n}_{S}(u)_{\ast} \Gamma^{n}_{\text{\rm ext}}ℱ{\quoi},
\]
两边均等同于
\[
\Gamma^{n}_{A}(F)^{\sim} \quad\mbox{\rm 当}\quad
ℱ=F^{\sim}{\quoi},
\]
\cqfd
\end{enumerate}

\setcounter{subsubsection}{34}
\subsubsection{}\etiquette[5.5.35]{XVII.subsubsec5.5.35}
处理 \hyperref[XVII.lem5.5.26]{5.5.26} 的情形 $2$. 用 \hyperref[XVII.lem5.5.22]{5.5.22} 的记号, 下图交换 (\hyperref[XVII.lem5.5.33]{5.5.33}):
$$\vcenter{\etiquette[5.5.35.1]{XVII.eq5.5.35.1}
\xymatrix@C=2cm{
  L\Gamma^{n} (R\Gamma(X,\ZZ/p)) \ar[r]^{k^n}_*+[o][F-]{\scriptstyle{1}}
  \ar[d]& R\Gamma(\Sym^{n}_{k}(X),\ZZ/p)\ar[d] \\
L\Gamma^{n}(R\Gamma(X,\Oo)) \ar[r]^{k^n}_*+[o][F-]{\scriptstyle{2}} & R\Gamma(\Sym^{n}_{k}(X),\Oo),
}}\leqno{(5.5.35.1)}
$$
并且\pageoriginale 由 \hyperref[XVII.prop5.5.34]{5.5.34}, 箭头 \ding{193} 为同构. $\Oo$ 的 Frobenius 自同态 $F:x\rightarrow x^{p}$ 在 $R\Gamma(X,\Oo)$ 及 $R\Gamma(\Sym^{n}_{k}(X),\Oo)$ 上 $p$-线性地作用. 因而它也在 $L\Gamma^{n}(R\Gamma(X,\Oo))$ 上 $p$-线性地作用. 可验证 $k^{n}$ 与 $F$ 可交换.

Artin--Schreier 理论给出正合列
$$
0\longrightarrow H^{i}(X,\ZZ/p)\longrightarrow
H^{i}(X,\Oo)\xrightarrow{F-I} H^{i}(X,\Oo)\longrightarrow
\sisi{\Oo}{0}{\quoi},
$$
对 $\Sym^{n}_{S}(X)$ 也类似. 要由 \ding{193} 为同构推出 \ding{192} 为同构, 只须验证以下引理.

\begin{enonce*}{引理 5.5.36}\etiquette[5.5.36]{XVII.lem5.5.36}
设 $k$ 为特征 $p>0$ 的完美域, $K\in C^{\geq
0}(\ZZ/p)$ 为有限维 $\ZZ/p$-向量空间复形, $L\in C^{\geq 0}(k)$ 为有限维 $k$-向量空间复形, $F$ 为 $L$ 的 $p$-线性自同态, 并设 $u:K\rightarrow L$ 为一个态射. 假设 $H^{\ast}(Fu)=H^{\ast}(u)$, 且序列
$$
0\longrightarrow H^{i}(K)\longrightarrow
H^{i}(L)\xrightarrow{F-I}H^{i}(L)\longrightarrow 0
$$
正合. 则\pageoriginale 序列
$$
0\longrightarrow H^{i}(\Gamma^{n}K)\longrightarrow
H^{i}(\Gamma^{n}L)\xrightarrow{F-I}H^{i}(\Gamma^{n}L)\longrightarrow 0
$$
正合.
\end{enonce*}

由于 $H^{\ast}(Fu)=H^{\ast}(u)$ 且 $\ZZ/p$ 为域, 存在同伦
$$
Fu-u=dH+Hd{\quoi}.
$$
由于域 $k$ 完美, 可以写成
$$
FH_{0}-H_{0}=H{\quoi}.
$$
用 $u-(dH_{0}+H_{0}d)$ 代替 $u$ 后, 可设 $Fu=u$.

函子 $S$ 将每个配有 $p$-线性自同态 $F$ 的 $k$-向量空间 $V$, 映为由满足 $Fv=v$ 的 $v\in
V$ 所成的 $\ZZ/p$-向量空间 $V^{S}$; 它是一个正合函子. 因而假设意味着态射
$$
K\longrightarrow L^{S}
$$
为拟同构. 只须证明函子 $S$ 与函子 $\Gamma^{n}$ 可交换. 这容易由下述众所周知的引理推出, 其证明留给读者:

\begin{enonce*}{引理 5.5.37}\etiquette[5.5.37]{XVII.lem5.5.37}
完美域 $k$ 上配有 $p$-线性自同态 $F$ 的有限维向量空间 $V$, 容许唯一的分解
$$
V=V_{0}+V_{n}
$$
满足 $V^{S}\otimes k\isomto  V_{0}$, $F(V_{n})\subset V_{n}$, 且 $F|V_{n}$ 幂零.
\end{enonce*}

这就完成了 \hyperref[XVII.thm5.5.21]{5.5.21} 的证明.


